%% ============================================================================
%% Chapter 20: API Lifecycle, Governance & Developer Experience
%% Mastering APIs: From Zero to Production Guru
%% ============================================================================
\chapter{API Lifecycle, Governance \& Developer Experience}
\label{ch:lifecycle}

%% ----------------------------------------------------------------------------
%% 1. Executive Summary
%% ----------------------------------------------------------------------------
Every API you have built in this book is immortal until proven otherwise.
The moment a consumer integrates against an endpoint --- the FastAPI
services of Chapter~\ref{ch:fastapi}, the versioned contracts of
Chapter~\ref{ch:versioning}, the rate-limited public surfaces of
Chapter~\ref{ch:rate_limiting} --- you have created an obligation that
outlives the enthusiasm that created it. Teams change, roadmaps pivot,
funding moves on, but the consumer's production workload keeps calling.
The difference between an API organization and a pile of endpoints is not
technical talent; it is \emph{lifecycle management}: knowing which states
an API can be in, what evidence is required to move between states, who is
allowed to decide, and how the people on the other side of the contract
experience every transition.

This chapter treats the API lifecycle as a \textbf{formal state machine}:
\textsf{design} $\rightarrow$ \textsf{review} $\rightarrow$
\textsf{published} (operate) $\rightarrow$ \textsf{deprecated}
$\rightarrow$ \textsf{sunset} $\rightarrow$ \textsf{retired}, with explicit
entry and exit criteria per state and a small set of legal transitions.
On top of that machine we build \textbf{design-first governance}: a style
guide that standardizes exactly the things consumers can observe (naming,
pagination, errors, authentication, versioning), enforced not by review
heroics but by linting --- Spectral-style rulesets executed as data in CI
--- plus a design-review process with service-level expectations, API
decision records, and the honesty to say that governance is a trade-off
between fleet-wide consistency and team autonomy. The resolution of that
trade-off is the \emph{paved path}: make the compliant way the easiest
way, and measure how many teams choose it voluntarily.

We then turn the lens around to \textbf{developer experience (DX)
engineering}. A developer portal with a searchable catalog, docs-as-code
pipelines that render OpenAPI into reference documentation
\cite{openapi31}, SDK generation (and a clear-eyed account of when
generated SDKs hurt more than they help), deterministic sandbox
environments, and above all \textbf{time-to-first-call (TTFC)} --- the
single activation metric that predicts whether an API will be adopted or
abandoned. Because APIs are \textbf{products}, we cover consumer research,
tiering and pricing models, support models, changelog discipline, and the
SLA/SLO/SLI relationship~\cite{beyer2016sre}, including the reasoning
behind the rule that a contractual SLA must be strictly looser than the
internal SLO that defends it. \textbf{Deprecation and sunsetting at scale}
gets a full treatment: discovering who actually calls what from gateway
telemetry (including the shadow APIs nobody owns), migration campaigns
with communication cadences, kill dates, brownouts, and escape hatches,
and the very different policies that apply to internal versus external
consumers. Finally we address \textbf{regulatory and data governance}
(PII boundaries in payloads and logs, retention, audit trails, GDPR-style
deletion propagation, classification labels in schemas) and
\textbf{platform economics} (build-versus-buy, funding models for platform
teams~\cite{newman2021microservices}, and measuring platform ROI in
adoption, TTFC, and incident reduction).

Everything is backed by five complete, offline, deterministic Python~3.11
tools against the synthetic Northwind Robotics universe: a mini
Spectral-style linter with rules-as-data and CI exit codes, an OpenAPI
$\rightarrow$ Markdown docs generator, a TTFC funnel analyzer over
synthetic onboarding events, a deprecation campaign tracker that prints
retirement readiness and the communication schedule, and a catalog
manifest generator that emits Backstage-style API entries from a folder of
specs.

\begin{keyconceptbox}
Governance that lives in documents is a wish; governance that lives in
code is a system. A style guide nobody can execute is enforced by the
scarcest resource you have --- senior reviewer attention --- and decays
the moment that attention is needed elsewhere. The platform-engineering
move is to convert every ``should'' into a lint rule, a template, a
generated artifact, or a default, so that compliance is the path of least
resistance and reviewer attention is reserved for the genuinely hard
questions: Is this the right API? Should it exist at all? What happens to
its consumers when it dies?
\end{keyconceptbox}

%% ----------------------------------------------------------------------------
\section{Learning Objectives \& Prerequisites}
\label{sec:ch20-objectives}

\subsection*{Learning Objectives}
After completing this chapter, its lab, and its exercises, you will be able to:

\begin{enumerate}[label=\textbf{LO-\arabic*.}, leftmargin=2.6cm]
  \item Model an API's life as a state machine with named states, legal
        transitions, and evidence-based entry/exit criteria, and identify
        which transition a given organizational event represents.
  \item Write an API style guide that standardizes consumer-observable
        surfaces (naming, pagination, errors, authentication, versioning),
        express it as an executable Spectral-style ruleset with custom
        rules, and wire it into CI with meaningful exit codes.
  \item Run a design-review process that respects team autonomy ---
        paved paths, API decision records, review SLAs --- and articulate
        precisely which decisions are centralized and which are not.
  \item Engineer developer experience as a measurable system: portals and
        catalogs, docs-as-code, SDK generation policy, sandboxes, and
        TTFC instrumentation with per-step funnel conversion.
  \item Operate an API as a product: tiering and pricing, support models,
        changelog discipline, and the SLA~$<$~SLO rule with its
        error-budget rationale~\cite{beyer2016sre}.
  \item Design and execute a deprecation campaign at scale: consumer
        discovery from gateway telemetry, communication cadence, brownouts
        as escape hatches, kill dates, and internal-versus-external policy
        differences (extending Chapter~\ref{ch:versioning}).
  \item Apply data-governance controls to API ecosystems: PII boundaries,
        retention, audit trails, deletion propagation, and schema-level
        classification labels.
  \item Evaluate platform economics: build-versus-buy for portals and
        gateways, platform team funding models, and an ROI scorecard in
        adoption, activation, and reliability terms.
\end{enumerate}

\subsection*{Prerequisites}
This chapter assumes you have completed Chapters~0--19 and are comfortable
with:

\begin{itemize}
  \item \textbf{Chapter~\ref{ch:fastapi}} (building APIs with FastAPI) ---
        the services being governed here are the ones you learned to
        build; \textbf{Chapter~\ref{ch:versioning}} (versioning \&
        evolution) --- deprecation and sunset policy extend the
        compatibility algebra developed there; \textbf{Chapter~\ref{ch:testing}}
        (testing \& quality engineering) --- lint gates are another
        quality gate in the same CI supply chain.
  \item OpenAPI~3.1 document structure~\cite{openapi31} from
        Chapters~\ref{ch:fastapi} and~\ref{ch:testing}: \texttt{info},
        \texttt{paths}, operations, \texttt{components.schemas}, and
        extension properties (\texttt{x-*}).
  \item RFC~9457 problem details (Chapter~\ref{ch:problem_details}),
        cursor pagination (Chapter~\ref{ch:pagination}), and
        rate-limit/deprecation headers (Chapters~\ref{ch:rate_limiting}
        and~\ref{ch:versioning}), since the style guide standardizes
        exactly these surfaces.
  \item SLI/SLO vocabulary and error budgets from
        Chapter~\ref{ch:observability}~\cite{beyer2016sre}.
  \item Python~3.11+, PyYAML, and pytest; all code runs offline on
        Windows/PowerShell with synthetic Northwind Robotics data only.
\end{itemize}

%% ----------------------------------------------------------------------------
\section{Deep Technical Mechanics \& Architectural Concepts}
\label{sec:ch20-mechanics}

\subsection{The API Lifecycle as a State Machine}
\label{sec:ch20-statemachine}

An ``API lifecycle diagram'' drawn as a circle of buzzwords is decoration.
A lifecycle you can \emph{operate} is a state machine: a finite set of
states, a small alphabet of events, legal transitions, and --- the part
everyone skips --- the \textbf{evidence required} to take each transition.

\begin{definition}[API Lifecycle State]
An \emph{API lifecycle state} is a tuple $(s, O_s, E_s, X_s)$ where $s$ is
the state name, $O_s$ the set of operational obligations in force while in
$s$ (e.g., SLA coverage, on-call, changelog duty), $E_s$ the entry
criteria (evidence that must exist to enter), and $X_s$ the exit criteria
(evidence required to leave). A \emph{transition} $s \xrightarrow{e} s'$
is an event $e$ fired by a named role and recorded in an audit log.
\end{definition}

The seven states we operate at Northwind Robotics are
\textsf{design}, \textsf{review}, \textsf{published},
\textsf{deprecated}, \textsf{sunset}, and \textsf{retired} --- with
``operate'' being the normal dwell mode of \textsf{published}. The states
are not a maturity gradient; an API in \textsf{deprecated} is not
``worse'' than one in \textsf{published}, it is \emph{under a different
contract}. Each state changes what consumers may assume, what the provider
must do, and what the platform automates.

\begin{figure}[htbp]
\centering
\begin{tikzpicture}[
  font=\small, >=stealth, node distance=0.9cm and 1.15cm,
  statebox/.style={rectangle, draw, rounded corners, align=center,
    minimum width=2.5cm, minimum height=0.95cm},
  terminal/.style={rectangle, draw, double, rounded corners,
    align=center, minimum width=2.5cm, minimum height=0.95cm},
  ev/.style={font=\footnotesize\itshape, text=packtgray}]
  \node[statebox] (design) {\textsf{design}\\ \footnotesize spec in repo, no traffic};
  \node[statebox, right=of design] (review) {\textsf{review}\\ \footnotesize lint + design review};
  \node[statebox, right=of review] (published) {\textsf{published}\\ \footnotesize catalog, SLA, on-call};
  \node[statebox, right=of published] (deprecated) {\textsf{deprecated}\\ \footnotesize frozen, \texttt{Sunset} header};
  \node[statebox, right=of deprecated] (sunset) {\textsf{sunset}\\ \footnotesize brownouts, 410};
  \node[terminal, right=of sunset] (retired) {\textsf{retired}\\ \footnotesize catalog entry removed};

  \draw[->] (design) -- node[ev, above] {submit for review} (review);
  \draw[->] (review) -- node[ev, above] {approved, lint clean} (published);
  \draw[->, bend left=28] (review) to node[ev, above] {changes requested} (design);
  \draw[->] (published) -- node[ev, above] {sunset announced} (deprecated);
  \draw[->] (deprecated) -- node[ev, above] {kill date reached} (sunset);
  \draw[->] (sunset) -- node[ev, above] {post-sunset audit clean} (retired);
  \draw[->, bend right=30] (deprecated) to
    node[ev, below] {exception: un-deprecate (rare, ADR required)} (published);
  \path (published.south) ++(0,-0.75) node[note]
    {dwell state: operate --- monitor SLOs, ship backward-compatible changes, honor changelog};
\end{tikzpicture}
\caption{The API lifecycle as a state machine. Each transition is an
audited event with named preconditions; \textsf{published} is the dwell
state where normal operation (and only backward-compatible evolution) happens.}
\label{fig:ch20-lifecycle}
\end{figure}

Table~\ref{tab:ch20-states} gives the entry/exit criteria. Two details
matter more than the rest. First, \textbf{entry into \textsf{published}
requires evidence, not opinion}: a clean lint run, a completed design
review, an owner in the catalog, and a declared SLO. Second,
\textbf{\textsf{deprecated} is a first-class operational state with real
machinery}: the contract is frozen (no new features, security fixes only),
every response carries a \texttt{Sunset} header
(Chapter~\ref{ch:versioning}), and the migration campaign
(Section~\ref{sec:ch20-deprecation}) is running against a public kill
date.

\begin{table}[htbp]
\centering\small
\begin{tabularx}{\textwidth}{l X X}
\toprule
\textbf{State} & \textbf{Entry criteria (evidence)} & \textbf{Exit criteria (evidence)} \\
\midrule
design & problem statement + ADR approved to proceed & spec in repo; lint clean; review packet complete \\
review & lint gate green in CI; owner identified & review board sign-off; open questions resolved or ADR'd \\
published (operate) & catalog entry, rendered docs, declared SLO, on-call rotation, sandbox live & announcement of deprecation date + replacement path \\
deprecated & deprecation ADR; consumer census from gateway telemetry; \texttt{Sunset}/\texttt{Deprecation} headers live & kill date reached; residual traffic below policy floor \\
sunset & brownout schedule executed; final warnings sent & post-sunset audit: 410 traffic $\approx$ 0 for one audit window \\
retired & audit clean; catalog entry archived; runbooks deleted & terminal state --- no exit \\
\bottomrule
\end{tabularx}
\caption{Lifecycle states with evidence-based entry/exit criteria.}
\label{tab:ch20-states}
\end{table}

The state machine earns its keep in three ways. It makes
\textbf{ownership answerable at any moment}: every state names an
accountable role (in \textsf{design} the proposing team; from
\textsf{published} onward the owning team of record in the catalog). It
makes \textbf{tooling targetable}: lint gates fire on the
\textsf{design}$\rightarrow$\textsf{review} transition, catalog rendering
on \textsf{review}$\rightarrow$\textsf{published}, telemetry diffing on
\textsf{published}$\rightarrow$\textsf{deprecated}. And it makes
\textbf{audits trivial}: ``which APIs entered \textsf{deprecated} without
a consumer census?'' is a log query, not an inquisition.

\subsection{Design-First Governance}
\label{sec:ch20-governance}

\subsubsection{What a Style Guide Standardizes --- and What It Leaves Alone}
A style guide that tries to standardize everything standardizes nothing,
because teams route around it. The discipline is to standardize only the
\textbf{consumer-observable seam} and leave internals to the teams. Our
canonical list, built up through the book:

\begin{table}[htbp]
\centering\small
\begin{tabularx}{\textwidth}{l X l}
\toprule
\textbf{Surface} & \textbf{Standard (Northwind Robotics)} & \textbf{Chapter} \\
\midrule
Path naming & kebab-case segments, plural nouns, no verbs & \ref{ch:rest_style} \\
Operation IDs & camelCase, unique, verb+noun (\texttt{listRobots}) & \ref{ch:fastapi} \\
Pagination & cursor-based, \texttt{limit} + \texttt{cursor} + \texttt{nextCursor} & \ref{ch:pagination} \\
Errors & RFC~9457 \texttt{application/problem+json} on every 4xx/5xx & \ref{ch:problem_details} \\
Authentication & OAuth2 client credentials; scopes named \texttt{domain:action} & \ref{ch:api_security} \\
Versioning & URI major version; \texttt{info.version} strict semver; \texttt{Sunset} header on deprecation & \ref{ch:versioning} \\
Idempotency & \texttt{Idempotency-Key} on all unsafe operations & \ref{ch:idempotency} \\
Rate limiting & standard \texttt{RateLimit}-* headers + 429 with \texttt{Retry-After} & \ref{ch:rate_limiting} \\
\bottomrule
\end{tabularx}
\caption{The consumer-observable surfaces a style guide must pin down.}
\label{tab:ch20-styleguide}
\end{table}

Everything not in the table --- language, framework, datastore, internal
module layout --- is deliberately ungoverned. This is the
\textbf{consistency-versus-autonomy trade-off} made explicit: consistency
buys fleet-wide learnability (a developer who knows one Northwind API
knows them all) and tooling leverage (one docs renderer, one SDK
generator, one linter); autonomy buys team speed and local fit. The
mistake is paying for consistency you cannot cash: mandating internal
choices produces compliance theater, while under-specifying the seam
produces fifty dialects of pagination.

\subsubsection{Spectral-Style Linting: Rules as Data}
A style guide becomes real when it compiles. Spectral (the open-source
OpenAPI linter this section's code imitates) expresses rules as
\emph{data}: a YAML ruleset where each rule names a target (a JSONPath-ish
expression over the document), a severity, and a function (truthy, pattern,
casing, schema\ldots). The anatomy generalizes:

\begin{itemize}
  \item \textbf{Rule identity and severity.} Stable IDs
        (\texttt{operation-id-camel-case}) so violations can be
        referenced, suppressed, and trended; severities mapped to CI
        semantics (\texttt{error} fails the build, \texttt{warn} does
        not --- yet).
  \item \textbf{Targets.} Where the rule looks: every operation, every
        path segment, a single pointer like \texttt{/info/contact}.
  \item \textbf{Functions.} The check itself: casing, regex match,
        required-field presence, or a domain predicate like ``every
        operation declares an \texttt{application/problem+json} error
        response''.
  \item \textbf{Messages with JSON paths.} A violation that cannot point
        at the offending node (\texttt{\$.paths["/RobotParts"].get})
        doubles the cost of fixing it.
\end{itemize}

Custom rules are where governance encodes \emph{your} history. The rule
\texttt{pagination-required-on-lists} exists because a 2019 incident
(Chapter~\ref{ch:pagination}'s war story) traced an outage to an
unpaginated collection endpoint. A ruleset is an institutional memory with
a \texttt{main()} function.

\subsubsection{Design Review: Process, Not Tribunal}
With lint absorbing the mechanical checks, human review is reserved for
judgment: Is the resource model right? Is this operation idempotent? Does
this API deserve to exist? A healthy review process has four properties:
(1)~\textbf{review SLA} --- feedback within two business days, or the
review is waived by default (a slow gate teaches teams to bypass gates);
(2)~\textbf{a written packet} --- the spec, the lint report, and a short
ADR for each non-obvious decision; (3)~\textbf{advisory-by-default
verdicts} with a short list of blocking conditions (security, data
classification, contract-breaking changes to published APIs);
(4)~\textbf{appeals} --- a route past the reviewer to an architecture
council, used rarely and publicly.

\subsubsection{API Decision Records}
An ADR for APIs is a one-page document: \emph{context} (forces and
constraints), \emph{decision} (what we will do), \emph{status}
(proposed/accepted/superseded), \emph{consequences} (what becomes easier
and harder), and \emph{alternatives considered}. Two ADR habits pay
outsized dividends in API programs: recording \textbf{rejected
alternatives} (they are the answers to next year's ``why not GraphQL?''
mail; see Chapter~\ref{ch:graphql}), and linking ADRs \textbf{from the
spec} via an \texttt{x-decisions} extension so the rationale travels with
the contract.

\subsubsection{Paved Paths (Golden Paths)}
The synthesis of style guide, linter, and review is the \textbf{paved
path}: a blessed, fully tooled route from idea to published API ---
repository template with a compliant OpenAPI skeleton, pre-wired CI with
the ruleset, generated docs, automatic catalog registration, a sandbox,
and dashboards from Chapter~\ref{ch:observability}. Teams may leave the
paved path, but then they own the paving. Figure~\ref{fig:ch20-pavedpath}
shows the flow; the essential property is that \emph{every arrow is
automation}, and the humans only appear where judgment is irreducible.

\begin{figure}[htbp]
\centering
\begin{tikzpicture}[
  font=\small, >=stealth, node distance=0.8cm and 1.0cm]
  \node[flowstep] (design) {\textbf{design}\\ OpenAPI spec from repo template\\
    \footnotesize ADR for non-obvious decisions};
  \node[flowstep, right=of design] (lint) {\textbf{lint (CI)}\\ ruleset as data\\
    \footnotesize exit 1 blocks merge};
  \node[flowstep, right=of lint] (review) {\textbf{review}\\ judgment only\\
    \footnotesize 2-day SLA, advisory default};
  \node[flowstep, right=of review] (publish) {\textbf{publish}\\ registry + gateway route\\
    \footnotesize SLO declared, on-call set};
  \node[flowstep, right=of publish] (portal) {\textbf{portal}\\ catalog, docs, sandbox, keys\\
    \footnotesize TTFC clock starts};
  \draw[->] (design) -- (lint);
  \draw[->] (lint) -- (review);
  \draw[->] (review) -- (publish);
  \draw[->] (publish) -- (portal);
  \draw[->, bend left=25, dashed] (lint) to node[note, above] {violations} (design);
  \draw[->, bend right=30, dashed] (review) to node[note, below] {changes requested} (design);
  \node[modcap, below=0.5cm of review] {every solid arrow is automation; dashed arrows are feedback};
\end{tikzpicture}
\caption{The paved path: design $\rightarrow$ lint $\rightarrow$ review
$\rightarrow$ publish $\rightarrow$ portal. Compliance is the default
because the compliant route is the one with tooling.}
\label{fig:ch20-pavedpath}
\end{figure}

\subsection{Developer Experience Engineering}
\label{sec:ch20-dx}

Developer experience is not a documentation gloss applied at launch; it is
an engineering discipline with instrumentation, funnels, and budgets. The
unit of analysis is the \emph{developer journey}: discover $\rightarrow$
evaluate $\rightarrow$ activate $\rightarrow$ integrate $\rightarrow$
scale. Most API programs instrument only the last two stages, which is why
they cannot explain their adoption numbers.

\subsubsection{Developer Portals and Catalogs}
The portal's foundation is a \textbf{catalog of record}: every API, its
owner, its lifecycle state, its spec, and its SLO --- Backstage's
\texttt{kind: API} entity being the reference shape. A catalog answers the
two questions that otherwise consume engineering weeks: \emph{does an API
for X already exist?} and \emph{who owns API Y?} (the war story in
Section~\ref{sec:ch20-lab} is what happens when the second question has no
answer). The catalog is also where the lifecycle state machine becomes
visible: filtering by state turns governance from a spreadsheet ritual
into a query. Catalog entries must be \emph{generated} from the spec repo
(Listing~20.7), never hand-maintained --- a hand-maintained
catalog is a rumor with a search box.

\subsubsection{Docs-as-Code}
Docs-as-code means the reference documentation is a \emph{build artifact
of the contract}: OpenAPI in, rendered reference out (Redoc, Stoplight
Elements, or the stdlib renderer in Listing~20.4), versioned
with the spec, deployed by the same pipeline. The invariant worth
tattooing: \textbf{humans edit the spec; machines edit the docs}. The
moment someone hot-fixes the published HTML, docs and contract diverge and
trust in both collapses (Pitfall~\ref{pit:docs-drift}). Guides and
tutorials remain hand-written, but they too live in the repo and are
reviewed like code.

\subsubsection{SDK Generation from OpenAPI}
The openapi-generator flow is seductive: one spec, $N$ languages, clients
regenerated in CI on every spec change. When it works, it is a superpower
--- typed models, retry hooks, and auth wiring for free. The caveats are
real: generated code is idiomatic in \emph{no} language (Java-flavored
Python and Python-flavored Java), generated clients lag generator bugs and
language evolution, and an unmaintained SDK is \emph{worse than none}
because it carries your logo into its own bit-rot. Our decision rule:
\textbf{generate SDKs when} the spec is stable, the consumer count justifies
maintenance, and regeneration is wired to CI with publish automation;
\textbf{ship a great OpenAPI document plus a 15-minute quickstart instead}
when the API is young, consumers are few, or the team cannot commit to
owning the SDK as a product. An interactive sandbox plus copy-paste
\texttt{httpx} snippets (Chapter~\ref{ch:consuming_python}) covers the
long tail at near-zero marginal cost.

\subsubsection{Sandbox Environments}
A sandbox is a deterministic, self-service environment with synthetic data
--- the same offline-first discipline this book applies to itself. Its
purpose is to compress evaluation time: a developer should be able to make
a realistic call within minutes, against data that never changes under
their feet (seeded fixtures, fixed clocks where relevant), with no sales
call and no approval queue. Sandbox keys are issued instantly and
self-service; production keys may carry vetting. A sandbox that requires a
ticket is not a sandbox; it is a waiting room with a URL.

\subsubsection{TTFC: The Activation Metric}
\begin{definition}[Time to First Call (TTFC)]
For developer $d$, $\mathrm{TTFC}(d) = t_d(\text{first successful 2xx}) -
t_d(\text{signup})$, where ``successful'' means the developer received a
\emph{useful} response, not merely a response. Reported fleet-wide as the
median: $\mathrm{TTFC} = \mathrm{median}_d\,\mathrm{TTFC}(d)$.
\end{definition}

We instrument TTFC as a funnel --- \textsf{signup} $\rightarrow$
\textsf{key issued} $\rightarrow$ \textsf{first call} $\rightarrow$
\textsf{first 200} --- because the aggregate number hides \emph{where}
developers stall. With per-step counts $N_i$ and conversions
$C_i = N_i / N_{i-1}$, the funnel localizes the friction: a weak
\textsf{signup}$\rightarrow$\textsf{key} step points at key-issuance
bureaucracy; weak \textsf{key}$\rightarrow$\textsf{first call} points at
documentation and quickstart quality; weak \textsf{first
call}$\rightarrow$\textsf{first 200} points at authentication friction and
error-message quality (this is where Chapter~\ref{ch:problem_details}'s
\texttt{detail} field pays rent). TTFC is \emph{the} activation metric
because it is the earliest signal that correlates with long-term adoption:
developers who succeed in the first hour almost never churn in the first
month. Every DX investment --- portal search, quickstarts, sandboxes, SDKs
--- is justified, and prioritized, by its modeled effect on one funnel
step. Listing~20.5 computes the funnel from synthetic event
logs.

\subsubsection{API Analytics for Product Decisions}
Beyond TTFC, the analytics that change roadmap decisions are:
\textbf{endpoint adoption} (which operations are used at all --- candidates
for deprecation are born here), \textbf{error mix per consumer} (a
consumer whose 4xx rate spikes after a release is a support case you can
open \emph{before} the email arrives), \textbf{version distribution}
(Chapter~\ref{ch:versioning}'s migration telemetry), and \textbf{top
consumers by revenue and by support cost} --- the ratio of which is the
most honest pricing input you own. The ethical boundary: analytics observe
\emph{traffic}, not payload content, unless the consumer contract says
otherwise (Section~\ref{sec:ch20-datagov}).

\subsection{API as a Product}
\label{sec:ch20-product}

\subsubsection{Consumer Research and Tiering}
API product management begins with consumer research that is boringly
concrete: interview ten consumers, read their integration code, and
count their workarounds. Pricing and packaging then follow the value
metric --- the unit that scales with the consumer's success (calls, seats,
robots managed). A canonical three-tier structure: \textbf{free} (sandbox
plus low-rate production, community support, best-effort availability ---
the adoption engine), \textbf{usage-based} (pay per unit of the value
metric, business-hours support, contractual SLA), and \textbf{enterprise}
(committed volumes, dedicated support channel, custom SLOs, contractual
sunset protection). The tier boundary that matters most technically is
rate limits (Chapter~\ref{ch:rate_limiting}): tiers must be enforceable
at the gateway, or the pricing page is fiction.

\subsubsection{SLA, SLO, SLI --- and Why SLA $<$ SLO}
Recall the stack from Chapter~\ref{ch:observability}~\cite{beyer2016sre}:
an \textbf{SLI} is a measurement (fraction of requests succeeding in
under 300\,ms), an \textbf{SLO} is an internal target on that SLI
(99.9\% monthly), and an \textbf{SLA} is the contractual promise with
penalties (99.5\%, service credits below). The invariant is
$\text{SLA} < \text{SLO}$, and the rationale is the error budget. If the
internal objective is tighter than the external promise, the gap
$1 - \text{SLO}$ versus $1 - \text{SLA}$ is a buffer in which the team can
\emph{detect, decide, and act} --- page humans, roll back, freeze deploys
--- before any contractual breach. With monthly SLI $p$ measured against
an SLA floor $p_{\text{SLA}}$, budget remaining is
$(p - p_{\text{SLA}})/(1 - p_{\text{SLA}})$; it must still be positive
when the SLO budget $(p - p_{\text{SLO}})/(1 - p_{\text{SLO}})$ is
exhausted, which holds iff $\text{SLA} < \text{SLO}$. An SLA equal to or
tighter than the SLO means customers discover your incidents before your
error-budget policy does. (The reversed mistake is
Pitfall~\ref{pit:sla}.)

\subsubsection{Support Models and Changelog Discipline}
Support scales by self-service: docs, a searchable issue history, status
page, and only then humans --- community forum for free tier, ticketed
queue for paid, named contacts for enterprise. Changelog discipline is the
cheapest trust instrument in the product: every consumer-visible change,
classified (added/changed/deprecated/removed/fixed/security), dated, and
linked from the \texttt{Deprecation} machinery of
Chapter~\ref{ch:versioning}. A changelog that lags the deploy train
teaches consumers to distrust every other signal you send.

\subsubsection{Developer Relations as a Feedback Loop}
DevRel closes the loop between consumers and governance: funnel analytics
and support tickets flow into style-guide and paved-path evolution;
governance changes flow out through changelog, portal, and migration
guides. Figure~\ref{fig:ch20-loop} shows the loop. The failure mode to
guard against is DevRel as \emph{megaphone} --- outbound only --- which
produces beautiful launches of APIs nobody can integrate.

\begin{figure}[htbp]
\centering
\begin{tikzpicture}[
  font=\small, >=stealth, node distance=1.5cm and 2.6cm]
  \node[flowstep] (govern) {\textbf{governance}\\ style guide, rulesets, ADRs};
  \node[flowstep, right=of govern] (platform) {\textbf{platform}\\ lint CI, registry, portal, sandbox};
  \node[flowstep, below=of platform] (consumers) {\textbf{consumers}\\ integrate, activate, scale};
  \node[flowstep, left=of consumers] (telemetry) {\textbf{telemetry}\\ gateway census, TTFC funnel,\\ error mix, version drift};
  \draw[->] (govern) -- node[note, above] {rules as data} (platform);
  \draw[->] (platform) -- node[note, right] {paved path} (consumers);
  \draw[->] (consumers) -- node[note, below] {observed behavior} (telemetry);
  \draw[->] (telemetry) -- node[note, left] {evidence $\rightarrow$ rule changes} (govern);
  \draw[->, dashed, bend left=18] (telemetry) to
    node[note, above, sloped] {\footnotesize product signals (adoption, churn risk)} (platform);
  \node[modcap, below=0.4cm of telemetry] {};
\end{tikzpicture}
\caption{The governance feedback loop: standards flow to consumers through
tooling; observed consumer behavior flows back as evidence that revises
the standards.}
\label{fig:ch20-loop}
\end{figure}

\subsection{Deprecation \& Sunsetting at Scale}
\label{sec:ch20-deprecation}

Chapter~\ref{ch:versioning} taught the mechanics of signaling change
(\texttt{Deprecation} and \texttt{Sunset} headers, version negotiation);
this section is about the \emph{campaign}: retiring a version that
hundreds of consumers depend on, without breaking their production or
your reputation.

\subsubsection{Consumer Discovery: Who Uses What?}
You cannot migrate consumers you cannot name. The discovery order is:
(1)~\textbf{gateway telemetry} --- join API keys/client IDs to consumer
identities and aggregate calls per (consumer, version, endpoint); this is
ground truth because packets do not lie; (2)~\textbf{the catalog of
record} --- the declared mapping, useful for owners and escalation
contacts; (3)~\textbf{contract registries} from
Chapter~\ref{ch:testing}. The gap between (1) and (2) is where
\textbf{shadow consumers} live: integrations nobody declared, often built
by departed employees, and --- worse --- \emph{shadow APIs}: provider
endpoints serving real traffic with no catalog entry at all. Both are
found by diffing observed traffic against declared registrations, which is
why consumer discovery is a telemetry problem before it is a process
problem.

\subsubsection{Migration Campaigns}
A campaign is a dated communication plan with a hard end state, run like a
release. The skeleton (implemented in Listing~20.6):
\textbf{T-180} announce (changelog, portal banner, email to registered
owners, \texttt{Deprecation} header on); \textbf{T-90/T-60} reminders to
non-migrated consumers, escalating to their engineering leads;
\textbf{T-30} first \textbf{brownout} --- a short, announced window where
the old version returns 410/404, the escape hatch that flushes out
consumers who never read email; \textbf{T-14} final warning with sponsor
escalation; \textbf{T-7} second, longer brownout; \textbf{T-0} permanent
sunset (HTTP 410 with a problem+json body pointing at the migration
guide); \textbf{T+7} post-sunset audit of residual traffic. Escape
hatches: a \emph{documented} extension process (time-boxed, signed by the
consumer's director, carrying an explicit new kill date) and a per-consumer
allowlist so a laggard can be carved out of a brownout without delaying
the campaign for everyone else.

\subsubsection{The Organizational Cost of Breaking Changes}
Breaking changes externalize cost onto every consumer simultaneously:
each must schedule, implement, test, and deploy a migration --- work with
zero feature value to them. The total cost of a v1$\rightarrow$v2 break
scales with consumer count, which is why the calculus of
Chapter~\ref{ch:versioning} (expand--migrate--contract, parallel change)
is not politeness but economics: every backward-compatible expansion you
ship instead of a break is $N$ consumer migrations you did not force.
Budget for the campaign itself too: provider-side migration tooling,
shims, and the support surge around brownouts are real engineering weeks.

\subsubsection{Internal versus External Policy}
\begin{table}[htbp]
\centering\small
\begin{tabularx}{\textwidth}{l X X}
\toprule
\textbf{Policy lever} & \textbf{Internal consumers} & \textbf{External consumers} \\
\midrule
Notice period & 60--90 days; sprint-plannable & 6--12 months; often contractual \\
Enforcement & platform can mandate via CI/deploy gates & persuasion, pricing, brownouts; contract terms govern \\
Escape hatch & CTO-approved exception with expiry & paid extended-support tier, time-boxed \\
Discovery & catalog + telemetry + org chart & telemetry + contracts; no org chart to consult \\
Sunset mechanics & hard kill dates, brownouts acceptable & brownouts announced per contract; 410 body with migration link mandatory \\
\bottomrule
\end{tabularx}
\caption{Deprecation policy: internal versus external consumers.}
\label{tab:ch20-internal-external}
\end{table}

\subsection{Regulatory \& Data Governance}
\label{sec:ch20-datagov}

API governance that ignores data governance is a compliance incident on a
timer. Five controls, each cheap at design time and brutal retrofitted:

\begin{enumerate}
  \item \textbf{PII boundaries in payloads and logs.} Classify every
        schema field; payloads carry only what the consumer's purpose
        requires (data minimization), and logs capture metadata ---
        never payload bodies by default. Chapter~\ref{ch:observability}'s
        structured logging must enforce redaction lists, and
        Chapter~\ref{ch:api_security}'s scopes should be
        classification-aware.
  \item \textbf{Classification labels in schemas.} Machine-readable
        labels such as \texttt{x-classification: pii|internal|public} on
        schema properties (used by our specs and surfaced by the docs
        generator in Listing~20.4) let linters fail builds that
        log, export, or widen access to classified fields.
  \item \textbf{Retention policies.} Every persisted representation has a
        retention clock and an owner; API audit logs, analytics events,
        and backups each get explicit, possibly different, windows.
  \item \textbf{Audit trails.} Lifecycle transitions (who moved this API
        to \textsf{deprecated}, on what evidence), admin operations, and
        data-access events are append-only, timestamped, and attributable.
  \item \textbf{Deletion propagation.} A GDPR-style erasure request is a
        \emph{workflow across the API ecosystem}, not a row delete:
        primary store, derived stores, analytics events, caches
        (Chapter~\ref{ch:idempotency}), backups (handled by expiry, not
        surgical edits), and downstream consumers you pushed data to via
        async APIs (Chapter~\ref{ch:async_apis}) --- each acknowledged and
        auditable. Designing erasure as an event with tombstones and
        fan-out receipts, rehearsed quarterly, beats discovering your
        propagation gaps during a regulator's deadline.
\end{enumerate}

\subsection{Platform Economics}
\label{sec:ch20-economics}

\textbf{Build versus buy.} Buy (or adopt managed) when the capability is
undifferentiated and mature --- gateways, key management, hosted docs.
Build when the capability encodes \emph{your} conventions --- the ruleset,
the paved-path templates, the catalog enrichment that knows your org
chart. The composite answer is usual: a managed gateway, Backstage-style
open-source portal core, and thin bespoke glue (like this chapter's
manifest generator) that stamps your conventions onto commodity
infrastructure. \textbf{Funding models.} Platform teams fail when funded
as ticket-taking cost centers; they succeed as \emph{product teams with
internal customers}~\cite{newman2021microservices} --- funded on a
multi-year horizon, with adoption as their success metric and voluntary
uptake as the forcing function (a platform that must mandate usage is a
platform whose paved path is not paved well enough). \textbf{Measuring
ROI.} The honest scorecard has three columns: \emph{adoption} (APIs
onboarded to the paved path, share of traffic behind the gateway),
\emph{activation} (TTFC median and funnel conversion, trended), and
\emph{risk removed} (lint-blocked non-compliant designs, incidents
attributed to governed versus ungoverned APIs, sunset campaigns completed
without escalation). Report it quarterly; a platform that cannot show its
value eventually meets someone who cannot see it.

%% ----------------------------------------------------------------------------
\section{Production-Ready Code Implementation}
\label{sec:ch20-code}

All listings live under \texttt{code\textbackslash{}ch20\textbackslash{}}
and together form a \emph{mini platform}: lint a spec, generate its docs,
publish its catalog entry, measure onboarding, and run a deprecation
campaign. Everything is offline, deterministic (the tracker takes
\texttt{--today} explicitly; the funnel reads a fixed synthetic fixture),
and uses the Northwind Robotics synthetic universe only. Install once:

\begin{minted}{powershell}
cd code\ch20
pip install -r requirements.txt
python -m pytest -q
\end{minted}

\subsection{Listing 20.1 --- The Ruleset as Data}
\label{sec:ch20-listing-ruleset}

The style guide of Table~\ref{tab:ch20-styleguide} compiles into a YAML
ruleset. Six rules cover contact ownership, semver versioning, path and
operation naming, RFC~9457 error coverage, and pagination on collection
endpoints. Adding a rule is a \emph{data} change, reviewed in a pull
request like any other --- governance as pull request, not as meeting.

\begin{minted}{yaml}
# Northwind Robotics API style ruleset (rules as data).
# Each rule: id, severity (error|warn), kind, and kind-specific params.
rules:
  - id: info-contact-required
    severity: error
    kind: require-field
    pointer: /info/contact
    message: info.contact must name an owning team or alias.

  - id: semver-info-version
    severity: error
    kind: match
    pointer: /info/version
    pattern: '^\d+\.\d+\.\d+$'
    message: info.version must be strict semver (MAJOR.MINOR.PATCH).

  - id: path-segment-kebab-case
    severity: error
    kind: casing
    target: path_segment
    style: kebab
    message: Path segments must be kebab-case.

  - id: operation-id-camel-case
    severity: error
    kind: casing
    target: operation_id
    style: camel
    message: operationId must be camelCase.

  - id: problem-json-errors
    severity: error
    kind: problem-json-errors
    message: >-
      Every operation must declare a 4xx/5xx or default response carrying
      application/problem+json (RFC 9457).

  - id: pagination-required-on-lists
    severity: warn
    kind: pagination-required
    params: [limit, cursor]
    message: >-
      Collection GET endpoints must accept limit and cursor query
      parameters.
\end{minted}

\subsection{Listing 20.2 --- \texttt{mini\_spectral.py}: A Spectral-Style Linter}
\label{sec:ch20-listing-linter}

The linter loads the spec and ruleset, evaluates each rule \emph{kind}
(\texttt{require-field}, \texttt{match}, \texttt{casing},
\texttt{problem-json-errors}, \texttt{pagination-required}), and reports
violations with JSON paths. The exit code is the CI contract: 0 clean, 1
error-severity violations, 2 usage or parse failure --- so a pipeline can
distinguish ``your API is non-compliant'' from ``the gate itself broke''.

\begin{minted}{python}
r"""mini_spectral.py -- a miniature Spectral-style OpenAPI linter.

Loads an OpenAPI YAML document and a ruleset expressed as data (YAML),
evaluates every rule, reports violations with JSON paths, and returns a
CI-friendly exit code:

    0 -- no error-severity violations
    1 -- at least one error-severity violation
    2 -- usage error, unreadable file, or malformed document/ruleset

Usage (PowerShell):
    python mini_spectral.py .\specs\northwind_robots.yaml --ruleset .\ruleset.yaml
"""

from __future__ import annotations

import argparse
import re
import sys
from dataclasses import dataclass
from pathlib import Path
from typing import Any, Callable

import yaml

HTTP_METHODS = frozenset(
    {"get", "put", "post", "delete", "options", "head", "patch", "trace"}
)

CASING_PATTERNS: dict[str, re.Pattern[str]] = {
    "kebab": re.compile(r"^[a-z0-9]+(-[a-z0-9]+)*$"),
    "snake": re.compile(r"^[a-z0-9]+(_[a-z0-9]+)*$"),
    "camel": re.compile(r"^[a-z][a-zA-Z0-9]*$"),
}

SEVERITIES = ("error", "warn")


@dataclass(frozen=True)
class Violation:
    """A single rule breach located by a JSON-path-ish string."""

    rule_id: str
    severity: str
    path: str
    message: str

    def render(self) -> str:
        return (
            f"[{self.severity.upper():5}] {self.rule_id} "
            f"at {self.path}: {self.message}"
        )


# --------------------------------------------------------------------------
# Document navigation helpers
# --------------------------------------------------------------------------

def resolve_pointer(doc: Any, pointer: str) -> Any:
    """Resolve a minimal JSON Pointer ('/info/contact'); None if absent."""
    node = doc
    for token in pointer.strip("/").split("/"):
        if not isinstance(node, dict) or token not in node:
            return None
        node = node[token]
    return node


def json_path(*parts: str) -> str:
    """Build a readable JSON path, quoting keys that contain slashes."""
    out = "$"
    for part in parts:
        if re.fullmatch(r"[A-Za-z0-9_-]+", part):
            out += f".{part}"
        else:
            out += f'["{part}"]'
    return out


def iter_operations(doc: dict[str, Any]):
    """Yield (json_path, path_key, method, operation) for every operation."""
    paths = doc.get("paths") or {}
    if not isinstance(paths, dict):
        return
    for path_key, path_item in paths.items():
        if not isinstance(path_item, dict):
            continue
        for method, operation in path_item.items():
            if method.lower() in HTTP_METHODS and isinstance(operation, dict):
                yield (
                    json_path("paths", path_key, method.lower()),
                    path_key,
                    method.lower(),
                    operation,
                )


# --------------------------------------------------------------------------
# Rule kinds -- each takes (doc, rule) and yields Violations
# --------------------------------------------------------------------------

def check_require_field(
    doc: dict[str, Any], rule: dict[str, Any]
) -> list[Violation]:
    pointer = str(rule["pointer"])
    if resolve_pointer(doc, pointer) is None:
        return [
            Violation(
                rule["id"],
                rule["severity"],
                "$" + pointer.replace("/", "."),
                rule["message"],
            )
        ]
    return []


def check_match(doc: dict[str, Any], rule: dict[str, Any]) -> list[Violation]:
    pointer = str(rule["pointer"])
    value = resolve_pointer(doc, pointer)
    pattern = re.compile(str(rule["pattern"]))
    if value is None or not pattern.fullmatch(str(value)):
        return [
            Violation(
                rule["id"],
                rule["severity"],
                "$" + pointer.replace("/", "."),
                f"{rule['message']} (found: {value!r})",
            )
        ]
    return []


def check_casing(doc: dict[str, Any], rule: dict[str, Any]) -> list[Violation]:
    style = str(rule["style"])
    pattern = CASING_PATTERNS[style]
    violations: list[Violation] = []

    if rule["target"] == "operation_id":
        for op_path, _pk, _m, operation in iter_operations(doc):
            op_id = operation.get("operationId")
            if op_id is None or not pattern.fullmatch(str(op_id)):
                violations.append(
                    Violation(
                        rule["id"],
                        rule["severity"],
                        f"{op_path}.operationId",
                        f"{rule['message']} (found: {op_id!r})",
                    )
                )
    elif rule["target"] == "path_segment":
        paths = doc.get("paths") or {}
        for path_key in paths:
            for segment in str(path_key).split("/"):
                if not segment or (segment.startswith("{") and segment.endswith("}")):
                    continue
                if not pattern.fullmatch(segment):
                    violations.append(
                        Violation(
                            rule["id"],
                            rule["severity"],
                            json_path("paths", path_key),
                            f"{rule['message']} (segment: {segment!r})",
                        )
                    )
    else:  # defensive: ruleset is data too, and data can be wrong
        raise ValueError(f"unknown casing target: {rule['target']!r}")
    return violations


def check_problem_json_errors(
    doc: dict[str, Any], rule: dict[str, Any]
) -> list[Violation]:
    violations = []
    for op_path, _pk, _m, operation in iter_operations(doc):
        responses = operation.get("responses") or {}
        ok = False
        for status, response in responses.items():
            code = str(status)
            if not (code.startswith(("4", "5")) or code == "default"):
                continue
            content = (response or {}).get("content") or {}
            if "application/problem+json" in content:
                ok = True
                break
        if not ok:
            violations.append(
                Violation(
                    rule["id"],
                    rule["severity"],
                    f"{op_path}.responses",
                    rule["message"],
                )
            )
    return violations


def check_pagination_required(
    doc: dict[str, Any], rule: dict[str, Any]
) -> list[Violation]:
    required_params = [str(p) for p in rule.get("params", ["limit", "cursor"])]
    violations = []
    for op_path, path_key, method, operation in iter_operations(doc):
        is_collection_get = method == "get" and not str(path_key).rstrip("/").endswith("}")
        if not is_collection_get:
            continue
        present = {
            str(p.get("name"))
            for p in operation.get("parameters") or []
            if isinstance(p, dict) and p.get("in") == "query"
        }
        missing = [p for p in required_params if p not in present]
        if missing:
            violations.append(
                Violation(
                    rule["id"],
                    rule["severity"],
                    f"{op_path}.parameters",
                    f"{rule['message']} (missing: {', '.join(missing)})",
                )
            )
    return violations


RULE_KINDS: dict[
    str, Callable[[dict[str, Any], dict[str, Any]], list[Violation]]
] = {
    "require-field": check_require_field,
    "match": check_match,
    "casing": check_casing,
    "problem-json-errors": check_problem_json_errors,
    "pagination-required": check_pagination_required,
}


# --------------------------------------------------------------------------
# Engine
# --------------------------------------------------------------------------

def load_yaml_file(path: Path) -> Any:
    with path.open("r", encoding="utf-8") as handle:
        return yaml.safe_load(handle)


def lint(spec: dict[str, Any], ruleset: dict[str, Any]) -> list[Violation]:
    """Evaluate every rule in the ruleset against the spec document."""
    if not isinstance(spec, dict) or "openapi" not in spec:
        raise ValueError("document does not look like an OpenAPI spec")
    rules = (ruleset or {}).get("rules") or []
    violations: list[Violation] = []
    for rule in rules:
        kind = rule.get("kind")
        handler = RULE_KINDS.get(str(kind))
        if handler is None:
            raise ValueError(f"unknown rule kind: {kind!r}")
        if rule.get("severity") not in SEVERITIES:
            raise ValueError(f"rule {rule.get('id')!r}: bad severity")
        violations.extend(handler(spec, rule))
    return violations


def main(argv: list[str] | None = None) -> int:
    parser = argparse.ArgumentParser(description="Mini Spectral-style linter.")
    parser.add_argument("spec", type=Path, help="OpenAPI YAML file")
    parser.add_argument(
        "--ruleset", type=Path, required=True, help="Ruleset YAML file"
    )
    args = parser.parse_args(argv)

    try:
        spec = load_yaml_file(args.spec)
        ruleset = load_yaml_file(args.ruleset)
    except (OSError, yaml.YAMLError) as exc:
        print(f"error: cannot load input: {exc}", file=sys.stderr)
        return 2

    try:
        violations = lint(spec, ruleset)
    except ValueError as exc:
        print(f"error: {exc}", file=sys.stderr)
        return 2

    for violation in violations:
        print(violation.render())

    errors = sum(1 for v in violations if v.severity == "error")
    warnings = sum(1 for v in violations if v.severity == "warn")
    print(f"\n{args.spec.name}: {errors} error(s), {warnings} warning(s)")
    return 1 if errors else 0


if __name__ == "__main__":
    raise SystemExit(main())
\end{minted}

The deliberately non-compliant legacy warehouse spec
(\texttt{fixtures\textbackslash{}noncompliant\_spec.yaml}) trips five of
the six rules:

\begin{minted}{yaml}
openapi: 3.1.0
info:
  title: Northwind Robotics Warehouse Legacy API
  version: v2-final   # NOT semver: violates semver-info-version
  description: Legacy warehouse endpoints kept alive for one partner.
  # no contact block: violates info-contact-required
paths:
  /RobotParts:        # PascalCase segment: violates path-segment-kebab-case
    get:
      operationId: list_parts   # snake_case: violates operation-id-camel-case
      summary: List all parts without pagination
      responses:
        '200':
          description: Every part in one giant payload
          content:
            application/json:
              schema:
                type: array
                items:
                  $ref: '#/components/schemas/Part'
        '500':
          description: Boom, plain text
          content:
            text/plain:
              schema:
                type: string
  /robot-parts/{partId}:
    delete:
      operationId: deletePart
      summary: Delete a part
      parameters:
        - name: partId
          in: path
          required: true
          schema:
            type: string
      responses:
        '204':
          description: Deleted
        '404':
          description: Not found, bare JSON
          content:
            application/json:
              schema:
                type: object
components:
  schemas:
    Part:
      type: object
      properties:
        sku:
          type: string
        bin:
          type: string
\end{minted}

Running the linter against it produces one violation per breach, each
pinpointed by JSON path, and exit code 1:

\begin{minted}{text}
[ERROR] info-contact-required at $.info.contact: info.contact must name an owning team or alias.
[ERROR] semver-info-version at $.info.version: info.version must be strict semver (MAJOR.MINOR.PATCH). (found: 'v2-final')
[ERROR] path-segment-kebab-case at $.paths["/RobotParts"]: Path segments must be kebab-case. (segment: 'RobotParts')
[ERROR] operation-id-camel-case at $.paths["/RobotParts"].get.operationId: operationId must be camelCase. (found: 'list_parts')
[ERROR] problem-json-errors at $.paths["/RobotParts"].get.responses: Every operation must declare a 4xx/5xx or default response carrying application/problem+json (RFC 9457).
[ERROR] problem-json-errors at $.paths["/robot-parts/{partId}"].delete.responses: Every operation must declare a 4xx/5xx or default response carrying application/problem+json (RFC 9457).
[WARN ] pagination-required-on-lists at $.paths["/RobotParts"].get.parameters: Collection GET endpoints must accept limit and cursor query parameters. (missing: limit, cursor)

noncompliant_spec.yaml: 6 error(s), 1 warning(s)
\end{minted}

The compliant fleet spec passes clean (exit 0):

\begin{minted}{text}

northwind_robots.yaml: 0 error(s), 0 warning(s)
\end{minted}

\subsection{Listing 20.3 --- Linter Tests, Including the Non-Compliant Spec}
\label{sec:ch20-listing-linter-tests}

The test suite treats the exit-code contract as first-class behavior: the
compliant spec exits 0, the non-compliant one exits 1 with every expected
rule ID in the output, a missing file exits 2, and warn-severity
violations provably do not fail the build --- the property that lets you
roll out new rules as warnings before flipping them to errors.

\begin{minted}{python}
"""Tests for mini_spectral.py, including a deliberately non-compliant spec."""

from __future__ import annotations

from pathlib import Path

import pytest
import yaml

import mini_spectral as ms

HERE = Path(__file__).parent
RULESET = HERE / "ruleset.yaml"
COMPLIANT = HERE / "specs" / "northwind_robots.yaml"
NONCOMPLIANT = HERE / "fixtures" / "noncompliant_spec.yaml"


def load(path: Path) -> dict:
    return yaml.safe_load(path.read_text(encoding="utf-8"))


# --------------------------------------------------------------------------
# End-to-end CLI behaviour (exit codes are the CI contract)
# --------------------------------------------------------------------------

def test_compliant_spec_exits_zero(capsys: pytest.CaptureFixture[str]) -> None:
    code = ms.main([str(COMPLIANT), "--ruleset", str(RULESET)])
    out = capsys.readouterr().out
    assert code == 0
    assert "0 error(s)" in out


def test_noncompliant_spec_exits_one_and_names_rules(
    capsys: pytest.CaptureFixture[str],
) -> None:
    code = ms.main([str(NONCOMPLIANT), "--ruleset", str(RULESET)])
    out = capsys.readouterr().out
    assert code == 1
    for rule_id in (
        "info-contact-required",
        "semver-info-version",
        "path-segment-kebab-case",
        "operation-id-camel-case",
        "problem-json-errors",
        "pagination-required-on-lists",
    ):
        assert rule_id in out
    # violations carry JSON paths
    assert '$.paths["/RobotParts"].get.operationId' in out
    assert "6 error(s), 1 warning(s)" in out


def test_missing_file_exits_two(capsys: pytest.CaptureFixture[str]) -> None:
    code = ms.main([str(HERE / "nope.yaml"), "--ruleset", str(RULESET)])
    assert code == 2
    assert "cannot load input" in capsys.readouterr().err


def test_not_an_openapi_doc_exits_two(
    tmp_path: Path, capsys: pytest.CaptureFixture[str]
) -> None:
    bogus = tmp_path / "bogus.yaml"
    bogus.write_text("just: a mapping\n", encoding="utf-8")
    code = ms.main([str(bogus), "--ruleset", str(RULESET)])
    assert code == 2
    assert "OpenAPI" in capsys.readouterr().err


# --------------------------------------------------------------------------
# Unit behaviour of individual rule kinds
# --------------------------------------------------------------------------

def test_casing_patterns() -> None:
    assert ms.CASING_PATTERNS["kebab"].fullmatch("robot-parts")
    assert not ms.CASING_PATTERNS["kebab"].fullmatch("RobotParts")
    assert ms.CASING_PATTERNS["camel"].fullmatch("listRobots")
    assert not ms.CASING_PATTERNS["camel"].fullmatch("list_parts")


def test_resolve_pointer_missing() -> None:
    assert ms.resolve_pointer({"info": {}}, "/info/contact") is None
    assert ms.resolve_pointer({"info": {"contact": {}}}, "/info/contact") == {}


def test_template_path_params_are_not_linted_as_segments() -> None:
    doc = load(COMPLIANT)
    ruleset = {
        "rules": [
            {
                "id": "path-segment-kebab-case",
                "severity": "error",
                "kind": "casing",
                "target": "path_segment",
                "style": "kebab",
                "message": "kebab only",
            }
        ]
    }
    assert ms.lint(doc, ruleset) == []


def test_unknown_rule_kind_is_a_ruleset_error() -> None:
    doc = load(COMPLIANT)
    with pytest.raises(ValueError, match="unknown rule kind"):
        ms.lint(doc, {"rules": [{"id": "x", "severity": "error",
                                 "kind": "telepathy"}]})


def test_warn_severity_does_not_fail_build(
    capsys: pytest.CaptureFixture[str],
) -> None:
    # The compliant spec trips the pagination rule if we demand a
    # non-existent parameter -- warnings must not change the exit code.
    ruleset_text = RULESET.read_text(encoding="utf-8").replace(
        "params: [limit, cursor]", "params: [limit, cursor, locale]"
    )
    violations = ms.lint(load(COMPLIANT), yaml.safe_load(ruleset_text))
    assert violations and all(v.severity == "warn" for v in violations)
\end{minted}

\subsection{Listing 20.4 --- \texttt{docs\_generator.py}: OpenAPI $\rightarrow$ Markdown}
\label{sec:ch20-listing-docs}

Docs-as-code in eighty lines of stdlib-plus-PyYAML: endpoints with
parameter and response tables, the first response example as a fenced JSON
block, and schema tables that surface \texttt{x-classification} labels so
data governance is visible in the published reference. Because the docs
are generated, they cannot drift from the spec without the pipeline
noticing.

\begin{minted}{python}
r"""docs_generator.py -- render an OpenAPI YAML spec as Markdown reference docs.

Stdlib + PyYAML only. Docs-as-code in miniature: the spec is the single
source of truth and the rendered reference is a build artifact, so docs
can never drift from the contract without the build noticing.

Usage (PowerShell):
    python docs_generator.py .\specs\northwind_robots.yaml
    python docs_generator.py .\specs\northwind_robots.yaml --out fleet-api.md
"""

from __future__ import annotations

import argparse
import json
import sys
from pathlib import Path
from typing import Any

import yaml

HTTP_METHODS = ("get", "put", "post", "delete", "patch", "options", "head")


def schema_type(schema: dict[str, Any]) -> str:
    """Human-readable type label, resolving local $refs to their name."""
    if "$ref" in schema:
        return str(schema["$ref"]).rsplit("/", 1)[-1]
    base = schema.get("type", "any")
    if base == "array":
        return f"array[{schema_type(schema.get('items') or {})}]"
    if "enum" in schema:
        return "enum(" + ", ".join(str(v) for v in schema["enum"]) + ")"
    if "format" in schema:
        return f"{base} ({schema['format']})"
    return str(base)


def render_parameters(operation: dict[str, Any]) -> list[str]:
    params = operation.get("parameters") or []
    if not params:
        return []
    lines = [
        "",
        "| Name | In | Required | Type |",
        "|---|---|---|---|",
    ]
    for param in params:
        schema = param.get("schema") or {}
        lines.append(
            f"| `{param.get('name', '?')}` | {param.get('in', '?')} | "
            f"{'yes' if param.get('required') else 'no'} | "
            f"{schema_type(schema)} |"
        )
    return lines


def render_responses(operation: dict[str, Any]) -> list[str]:
    responses = operation.get("responses") or {}
    lines = ["", "| Status | Meaning | Body |", "|---|---|---|"]
    for status, response in responses.items():
        content = (response or {}).get("content") or {}
        bodies = ", ".join(
            f"{media} -> {schema_type((cfg or {}).get('schema') or {})}"
            for media, cfg in content.items()
        ) or "-"
        lines.append(
            f"| `{status}` | {(response or {}).get('description', '')} "
            f"| {bodies} |"
        )
    return lines


def render_example(operation: dict[str, Any]) -> list[str]:
    """Emit the first response example found, as a fenced JSON block."""
    for _status, response in (operation.get("responses") or {}).items():
        for _media, cfg in ((response or {}).get("content") or {}).items():
            if isinstance(cfg, dict) and "example" in cfg:
                return [
                    "",
                    "Example:",
                    "",
                    "```json",
                    json.dumps(cfg["example"], indent=2),
                    "```",
                ]
    return []


def render_markdown(spec: dict[str, Any]) -> str:
    info = spec.get("info") or {}
    lines: list[str] = [
        f"# {info.get('title', 'Untitled API')}",
        "",
        f"**Version:** {info.get('version', 'n/a')}",
        "",
        (info.get("description") or "").strip(),
        "",
    ]
    for server in spec.get("servers") or []:
        lines.append(f"- Server: `{server.get('url')}` "
                     f"({server.get('description', '')})")
    lines.append("")
    lines.append("## Endpoints")

    paths = spec.get("paths") or {}
    for path_key, path_item in paths.items():
        for method in HTTP_METHODS:
            operation = (path_item or {}).get(method)
            if not isinstance(operation, dict):
                continue
            lines += [
                "",
                f"### `{method.upper()} {path_key}`",
                "",
                operation.get("summary", ""),
            ]
            lines += render_parameters(operation)
            lines += render_responses(operation)
            lines += render_example(operation)
            lines.append("")

    schemas = ((spec.get("components") or {}).get("schemas")) or {}
    if schemas:
        lines.append("## Schemas")
    for name, schema in schemas.items():
        required = set(schema.get("required") or [])
        lines += [
            "",
            f"### {name}",
            "",
            "| Property | Type | Required | Notes |",
            "|---|---|---|---|",
        ]
        for prop, prop_schema in (schema.get("properties") or {}).items():
            notes = prop_schema.get("description", "")
            if "x-classification" in prop_schema:
                notes = (notes + " " if notes else "") + (
                    f"[classification: {prop_schema['x-classification']}]"
                )
            lines.append(
                f"| `{prop}` | {schema_type(prop_schema)} | "
                f"{'yes' if prop in required else 'no'} | {notes} |"
            )
    lines.append("")
    return "\n".join(lines)


def main(argv: list[str] | None = None) -> int:
    parser = argparse.ArgumentParser(description="OpenAPI -> Markdown docs.")
    parser.add_argument("spec", type=Path, help="OpenAPI YAML file")
    parser.add_argument("--out", type=Path, default=None,
                        help="Write Markdown here (default: stdout)")
    args = parser.parse_args(argv)

    try:
        spec = yaml.safe_load(args.spec.read_text(encoding="utf-8"))
    except (OSError, yaml.YAMLError) as exc:
        print(f"error: cannot load spec: {exc}", file=sys.stderr)
        return 2
    if not isinstance(spec, dict) or "openapi" not in spec:
        print("error: not an OpenAPI document", file=sys.stderr)
        return 2

    markdown = render_markdown(spec)
    if args.out:
        args.out.write_text(markdown, encoding="utf-8")
        print(f"wrote {args.out} ({len(markdown)} chars)")
    else:
        print(markdown)
    return 0


if __name__ == "__main__":
    raise SystemExit(main())
\end{minted}

Demo output (truncated to the first endpoint; the full document covers all
operations and schemas):

\begin{minted}{text}
# Northwind Robotics Fleet API

**Version:** 2.4.0

Manages the Northwind Robotics synthetic robot fleet: registration, health telemetry summaries, and decommissioning workflows.

- Server: `https://api.northwind-robotics.example/fleet/v2` (Production (synthetic))

## Endpoints

### `GET /robots`

List robots in the fleet

| Name | In | Required | Type |
|---|---|---|---|
| `limit` | query | no | integer |
| `cursor` | query | no | string |

| Status | Meaning | Body |
|---|---|---|
| `200` | A page of robots | application/json -> RobotPage |
| `default` | Unexpected error | application/problem+json -> Problem |

Example:

```json
{
  "items": [
    {
(...)
\end{minted}

\subsection{Listing 20.5 --- \texttt{ttfc\_funnel.py}: Activation Analytics}
\label{sec:ch20-listing-ttfc}

The funnel analyzer reads a JSONL event log (one synthetic onboarding
event per line), reconstructs per-developer journeys, and reports the
ordered funnel with per-step and cumulative conversion plus median
inter-step gaps. It is defensive by design: malformed lines are counted
and skipped, non-funnel events (like \texttt{docs\_viewed}) are ignored,
and a developer is only counted at a step if all prior steps are present
--- out-of-order telemetry cannot fabricate conversions. The fixture (head
shown below) holds 222 events for 60 synthetic developers:

\begin{minted}{text}
{"ts": "2026-03-02T09:07:00Z", "developer_id": "dev-001", "event": "signup", "channel": "docs-site"}
{"ts": "2026-03-02T09:09:00Z", "developer_id": "dev-001", "event": "docs_viewed", "page": "quickstart"}
{"ts": "2026-03-02T09:14:00Z", "developer_id": "dev-002", "event": "signup", "channel": "docs-site"}
{"ts": "2026-03-02T09:16:00Z", "developer_id": "dev-002", "event": "docs_viewed", "page": "quickstart"}
{"ts": "2026-03-02T09:21:00Z", "developer_id": "dev-003", "event": "signup", "channel": "conference"}
{"ts": "2026-03-02T09:23:00Z", "developer_id": "dev-003", "event": "docs_viewed", "page": "quickstart"}
{"ts": "2026-03-02T09:28:00Z", "developer_id": "dev-004", "event": "signup", "channel": "docs-site"}
{"ts": "2026-03-02T09:30:00Z", "developer_id": "dev-004", "event": "docs_viewed", "page": "quickstart"}
(...)
\end{minted}

\begin{minted}{python}
r"""ttfc_funnel.py -- onboarding funnel and time-to-first-call analytics.

Reads a JSONL event log (one synthetic onboarding event per line) and
computes the activation funnel

    signup -> key_issued -> first_call -> first_200

with per-step absolute counts, step-to-step conversion, and median
elapsed minutes between steps. TTFC is reported as the median time from
signup to first_200 -- the first moment a developer received a *useful*
response.

Usage (PowerShell):
    python ttfc_funnel.py .\fixtures\onboarding_events.jsonl
"""

from __future__ import annotations

import argparse
import json
import statistics
import sys
from dataclasses import dataclass, field
from datetime import datetime
from pathlib import Path

FUNNEL_STEPS = ("signup", "key_issued", "first_call", "first_200")


@dataclass
class DeveloperJourney:
    """First occurrence of each funnel step for one developer."""

    first_seen: dict[str, datetime] = field(default_factory=dict)


def parse_ts(raw: str) -> datetime:
    return datetime.strptime(raw, "%Y-%m-%dT%H:%M:%SZ")


def load_journeys(path: Path) -> tuple[dict[str, DeveloperJourney], int]:
    """Build per-developer journeys; returns (journeys, skipped_lines)."""
    journeys: dict[str, DeveloperJourney] = {}
    skipped = 0
    with path.open("r", encoding="utf-8") as handle:
        for line in handle:
            line = line.strip()
            if not line:
                continue
            try:
                record = json.loads(line)
                dev = str(record["developer_id"])
                event = str(record["event"])
                ts = parse_ts(str(record["ts"]))
            except (ValueError, KeyError, TypeError):
                skipped += 1  # never let one bad line kill the report
                continue
            if event not in FUNNEL_STEPS:
                continue  # telemetry like docs_viewed is not a funnel step
            journey = journeys.setdefault(dev, DeveloperJourney())
            journey.first_seen.setdefault(event, ts)
    return journeys, skipped


def reached(journeys: dict[str, DeveloperJourney], step: str) -> set[str]:
    """Developers who completed this step AND every step before it."""
    index = FUNNEL_STEPS.index(step)
    required = FUNNEL_STEPS[: index + 1]
    return {
        dev
        for dev, journey in journeys.items()
        if all(s in journey.first_seen for s in required)
    }


def median_gap_minutes(
    journeys: dict[str, DeveloperJourney], start: str, end: str
) -> float | None:
    gaps = [
        (j.first_seen[end] - j.first_seen[start]).total_seconds() / 60.0
        for j in journeys.values()
        if start in j.first_seen and end in j.first_seen
    ]
    return statistics.median(gaps) if gaps else None


def build_report(journeys: dict[str, DeveloperJourney]) -> list[str]:
    lines = [
        "Onboarding funnel (unique developers, ordered steps)",
        "=" * 68,
        f"{'step':<14}{'reached':>9}{'step conv.':>12}{'cum. conv.':>12}"
        f"{'median gap (min)':>18}",
        "-" * 68,
    ]
    total = len(reached(journeys, FUNNEL_STEPS[0])) or 1  # avoid div-by-zero
    previous_count: int | None = None
    previous_step: str | None = None
    for step_name in FUNNEL_STEPS:
        count = len(reached(journeys, step_name))
        step_conv = (
            100.0 * count / previous_count
            if previous_count
            else 100.0
        )
        gap = (
            median_gap_minutes(journeys, previous_step, step_name)
            if previous_step
            else None
        )
        lines.append(
            f"{step_name:<14}{count:>9}{step_conv:>11.1f}%"
            f"{100.0 * count / total:>11.1f}%"
            f"{('n/a' if gap is None else f'{gap:.0f}'):>18}"
        )
        previous_count, previous_step = count, step_name

    ttfc = median_gap_minutes(journeys, "signup", "first_200")
    lines += [
        "-" * 68,
        f"TTFC (median signup -> first_200): "
        f"{('n/a' if ttfc is None else f'{ttfc:.0f} minutes')} "
        f"({('' if ttfc is None else f'{ttfc / 60:.1f} h')})",
    ]
    return lines


def main(argv: list[str] | None = None) -> int:
    parser = argparse.ArgumentParser(description="TTFC funnel analytics.")
    parser.add_argument("events", type=Path, help="JSONL onboarding events")
    args = parser.parse_args(argv)

    try:
        journeys, skipped = load_journeys(args.events)
    except OSError as exc:
        print(f"error: cannot read events: {exc}", file=sys.stderr)
        return 2
    if not journeys:
        print("error: no usable events found", file=sys.stderr)
        return 2

    for line in build_report(journeys):
        print(line)
    if skipped:
        print(f"note: skipped {skipped} malformed line(s)", file=sys.stderr)
    return 0


if __name__ == "__main__":
    raise SystemExit(main())
\end{minted}

The printed funnel table is the artifact you put in front of the DX
review board:

\begin{minted}{text}
Onboarding funnel (unique developers, ordered steps)
====================================================================
step            reached  step conv.  cum. conv.  median gap (min)
--------------------------------------------------------------------
signup               60      100.0%      100.0%               n/a
key_issued           45       75.0%       75.0%                69
first_call           33       73.3%       55.0%               194
first_200            24       72.7%       40.0%                93
--------------------------------------------------------------------
TTFC (median signup -> first_200): 431 minutes (7.2 h)
\end{minted}

Read it as a product manager: 75\% of signups obtain a key in a median 69
minutes, but only 73\% of those make a first call, and the
key$\rightarrow$call gap (194 minutes median) is the largest in the funnel
--- the quickstart, not key issuance, is where Northwind Robotics loses
developers. One instrumentation subtlety: \textsf{first call} and
\textsf{first 200} are separate steps because an 11\% slice of first calls
in this fixture return 401 (expired sandbox keys), which is exactly the
kind of friction TTFC exists to expose.

\subsection{Listing 20.6 --- \texttt{deprecation\_tracker.py}: Campaign Management}
\label{sec:ch20-listing-deprecation}

The tracker models a campaign as a registry --- API, version, replacement,
announce and kill dates, and a consumer census with per-consumer status,
traffic, and notes --- and renders two artifacts: the retirement-readiness
report (residual traffic, blockers, verdict) and the full communication
schedule derived from the kill date. The exit code is 1 while any
non-migrated consumer still has traffic, so the retirement runbook can be
gated in CI. The registry (note the last entry --- a shadow integration
found by gateway telemetry, with no owner in the catalog):

\begin{minted}{yaml}
# Synthetic deprecation campaign registry for the legacy Fleet API v1.
api: fleet
version: v1
replacement: v2
announce_date: "2026-03-02"
kill_date: "2026-09-01"   # hard sunset: version returns HTTP 410 after this
consumers:
  - name: fleet-dashboard
    owner: team-telemetry
    tier: internal
    status: migrated
    weekly_calls: 0
    last_contact: "2026-05-04"
    notes: Cut over to v2 in sprint 12.
  - name: maintenance-planner
    owner: team-operations
    tier: internal
    status: in_progress
    weekly_calls: 8200
    last_contact: "2026-06-29"
    notes: Blocked on v2 bulk endpoint; ETA 2026-08-01.
  - name: partner-robologix
    owner: external:robologix
    tier: external
    status: in_progress
    weekly_calls: 41000
    last_contact: "2026-06-21"
    notes: Contracted migration window ends 2026-08-15.
  - name: warehouse-scanner
    owner: team-warehouse
    tier: internal
    status: not_started
    weekly_calls: 1300
    last_contact: "2026-04-10"
    notes: Team dissolved; ownership transferred to team-fleet-platform.
  - name: legacy-billing-hook
    owner: unknown
    tier: internal
    status: unresponsive
    weekly_calls: 95
    last_contact: null
    notes: Discovered via gateway telemetry; no owner in catalog. Shadow API!
\end{minted}

\begin{minted}{python}
r"""deprecation_tracker.py -- retirement readiness + communication schedule.

Reads a deprecation campaign registry (YAML) and produces:

  1. a retirement-readiness report (per consumer, plus traffic-weighted
     migration percentage), and
  2. the communication schedule relative to the kill date
     (announcement, reminders, brownouts, sunset, post-sunset audit).

The report is deterministic: the evaluation date is supplied explicitly
(--today) so CI output and this book's examples never drift.

Usage (PowerShell):
    python deprecation_tracker.py .\fixtures\deprecation_registry.yaml --today 2026-07-15
"""

from __future__ import annotations

import argparse
import sys
from dataclasses import dataclass
from datetime import date, timedelta
from pathlib import Path
from typing import Any

import yaml

# Statuses that block retirement even at zero reported traffic.
ACTIVE_STATUSES = ("in_progress", "not_started")

# Communication cadence, as offsets (in days) from the kill date.
# Brownouts (scheduled short 410 windows) are the escape hatch that flushes
# out consumers who ignore email -- see Chapters 9 and 14.
CADENCE: tuple[tuple[int, str, str], ...] = (
    (-180, "announcement", "all consumers: email + changelog + portal banner"),
    (-90, "reminder", "all consumers not yet migrated"),
    (-60, "reminder", "non-migrated consumers + their owning teams' leads"),
    (-30, "brownout #1", "1-hour 410 window, business hours, with notice"),
    (-14, "final warning", "non-migrated consumers, escalation to sponsors"),
    (-7, "brownout #2", "4-hour 410 window with notice"),
    (0, "sunset", "permanent HTTP 410 + final changelog entry"),
    (7, "post-sunset audit", "gateway log sweep for residual 410 traffic"),
)


@dataclass(frozen=True)
class Consumer:
    name: str
    owner: str
    tier: str
    status: str
    weekly_calls: int
    notes: str = ""


def load_registry(path: Path) -> dict[str, Any]:
    data = yaml.safe_load(path.read_text(encoding="utf-8"))
    if not isinstance(data, dict) or "consumers" not in data:
        raise ValueError("registry must contain a 'consumers' list")
    return data


def parse_consumers(registry: dict[str, Any]) -> list[Consumer]:
    consumers = []
    for raw in registry["consumers"]:
        consumers.append(
            Consumer(
                name=str(raw["name"]),
                owner=str(raw.get("owner", "unknown")),
                tier=str(raw.get("tier", "internal")),
                status=str(raw.get("status", "not_started")),
                weekly_calls=int(raw.get("weekly_calls", 0)),
                notes=str(raw.get("notes", "")),
            )
        )
    return consumers


def readiness_report(
    registry: dict[str, Any], consumers: list[Consumer], today: date
) -> list[str]:
    kill = date.fromisoformat(str(registry["kill_date"]))
    # Migrated consumers no longer call v1, so the meaningful readiness
    # metric is residual traffic: calls still arriving from consumers that
    # have NOT finished migrating.
    residual = [c for c in consumers if c.status != "migrated"]
    residual_calls = sum(c.weekly_calls for c in residual)

    lines = [
        f"Retirement readiness: {registry['api']} {registry['version']} "
        f"(replacement: {registry.get('replacement', 'n/a')})",
        "=" * 72,
        f"today: {today.isoformat()}   kill date: {kill.isoformat()}   "
        f"days remaining: {(kill - today).days}",
        "",
        f"{'consumer':<22}{'tier':<10}{'status':<14}{'calls/wk':>10}  owner",
        "-" * 72,
    ]
    for c in sorted(consumers, key=lambda c: -c.weekly_calls):
        lines.append(
            f"{c.name:<22}{c.tier:<10}{c.status:<14}{c.weekly_calls:>10}  "
            f"{c.owner}"
        )
    lines += [
        "-" * 72,
        f"residual v1 traffic: {residual_calls} calls/wk across "
        f"{len(residual)} non-migrated consumer(s)",
        "",
    ]

    blockers = [
        c for c in consumers
        if c.status in ACTIVE_STATUSES
        or (c.status == "unresponsive" and c.weekly_calls > 0)
    ]
    if blockers:
        lines.append("verdict: NOT READY TO RETIRE -- blockers:")
        for c in blockers:
            lines.append(f"  - {c.name} ({c.status}, {c.weekly_calls} "
                         f"calls/wk): {c.notes}")
    else:
        lines.append("verdict: READY TO RETIRE -- no live traffic on "
                     "non-migrated consumers.")
    return lines


def schedule_report(registry: dict[str, Any]) -> list[str]:
    kill = date.fromisoformat(str(registry["kill_date"]))
    lines = [
        "",
        f"Communication schedule (kill date {kill.isoformat()})",
        "-" * 72,
        f"{'date':<12}{'T-minus':>9}  action / audience",
        "-" * 72,
    ]
    for offset, action, audience in CADENCE:
        when = kill + timedelta(days=offset)
        label = f"T{offset:+d}" if offset else "T-0"
        lines.append(f"{when.isoformat():<12}{label:>9}  {action}: {audience}")
    return lines


def main(argv: list[str] | None = None) -> int:
    parser = argparse.ArgumentParser(description="Deprecation campaign tracker.")
    parser.add_argument("registry", type=Path, help="campaign registry YAML")
    parser.add_argument(
        "--today",
        required=True,
        help="evaluation date, ISO format (kept explicit for determinism)",
    )
    args = parser.parse_args(argv)

    try:
        today = date.fromisoformat(args.today)
    except ValueError:
        print("error: --today must be ISO format, e.g. 2026-07-15",
              file=sys.stderr)
        return 2
    try:
        registry = load_registry(args.registry)
        consumers = parse_consumers(registry)
    except (OSError, ValueError, KeyError, yaml.YAMLError) as exc:
        print(f"error: cannot parse registry: {exc}", file=sys.stderr)
        return 2

    for line in readiness_report(registry, consumers, today):
        print(line)
    for line in schedule_report(registry):
        print(line)
    # Non-zero exit while blocked so CI can gate the retirement runbook.
    blocked = any(
        c.status in ACTIVE_STATUSES
        or (c.status == "unresponsive" and c.weekly_calls > 0)
        for c in consumers
    )
    return 1 if blocked else 0


if __name__ == "__main__":
    raise SystemExit(main())
\end{minted}

\begin{minted}{text}
Retirement readiness: fleet v1 (replacement: v2)
========================================================================
today: 2026-07-15   kill date: 2026-09-01   days remaining: 48

consumer              tier      status          calls/wk  owner
------------------------------------------------------------------------
partner-robologix     external  in_progress        41000  external:robologix
maintenance-planner   internal  in_progress         8200  team-operations
warehouse-scanner     internal  not_started         1300  team-warehouse
legacy-billing-hook   internal  unresponsive          95  unknown
fleet-dashboard       internal  migrated               0  team-telemetry
------------------------------------------------------------------------
residual v1 traffic: 50595 calls/wk across 4 non-migrated consumer(s)

verdict: NOT READY TO RETIRE -- blockers:
  - maintenance-planner (in_progress, 8200 calls/wk): Blocked on v2 bulk endpoint; ETA 2026-08-01.
  - partner-robologix (in_progress, 41000 calls/wk): Contracted migration window ends 2026-08-15.
  - warehouse-scanner (not_started, 1300 calls/wk): Team dissolved; ownership transferred to team-fleet-platform.
  - legacy-billing-hook (unresponsive, 95 calls/wk): Discovered via gateway telemetry; no owner in catalog. Shadow API!

Communication schedule (kill date 2026-09-01)
------------------------------------------------------------------------
date          T-minus  action / audience
------------------------------------------------------------------------
2026-03-05      T-180  announcement: all consumers: email + changelog + portal banner
2026-06-03       T-90  reminder: all consumers not yet migrated
2026-07-03       T-60  reminder: non-migrated consumers + their owning teams' leads
2026-08-02       T-30  brownout #1: 1-hour 410 window, business hours, with notice
2026-08-18       T-14  final warning: non-migrated consumers, escalation to sponsors
2026-08-25        T-7  brownout #2: 4-hour 410 window with notice
2026-09-01        T-0  sunset: permanent HTTP 410 + final changelog entry
2026-09-08        T+7  post-sunset audit: gateway log sweep for residual 410 traffic
\end{minted}

\subsection{Listing 20.7 --- \texttt{catalog\_manifest.py}: Backstage-Style Entries}
\label{sec:ch20-listing-catalog}

The manifest generator scans a directory of OpenAPI specs and emits one
catalog document per API, pulling lifecycle state and owner from the
spec's \texttt{info.x-lifecycle} block so the catalog is generated from
the contract of record rather than curated by hand. Non-OpenAPI files in
the folder are skipped with a warning, not a crash.

\begin{minted}{python}
r"""catalog_manifest.py -- generate Backstage-style API catalog entries.

Scans a directory of OpenAPI YAML specs and emits one catalog manifest
document per spec (apiVersion backstage.io/v1alpha1, kind: API), with
ownership and lifecycle read from the spec's info.x-lifecycle block so
the catalog, like the docs, is generated from the contract rather than
maintained by hand.

Usage (PowerShell):
    python catalog_manifest.py .\specs
    python catalog_manifest.py .\specs --out catalog-apis.yaml
"""

from __future__ import annotations

import argparse
import re
import sys
from pathlib import Path
from typing import Any

import yaml

DEFAULT_OWNER = "platform-team"
DEFAULT_LIFECYCLE = "production"


def slugify(title: str) -> str:
    slug = re.sub(r"[^a-z0-9]+", "-", title.lower()).strip("-")
    return slug or "unnamed-api"


def manifest_for(spec_path: Path, spec: dict[str, Any], root: Path) -> dict[str, Any]:
    info = spec.get("info") or {}
    lifecycle_block = info.get("x-lifecycle") or {}
    rel_path = spec_path.relative_to(root.parent).as_posix()
    return {
        "apiVersion": "backstage.io/v1alpha1",
        "kind": "API",
        "metadata": {
            "name": slugify(str(info.get("title", spec_path.stem))),
            "description": " ".join(
                str(info.get("description", "")).split()
            )[:200],
            "annotations": {
                "northwind/spec-version": str(info.get("version", "0.0.0")),
            },
        },
        "spec": {
            "type": "openapi",
            "lifecycle": str(lifecycle_block.get("state", DEFAULT_LIFECYCLE)),
            "owner": str(lifecycle_block.get("owner", DEFAULT_OWNER)),
            "definition": {"$text": f"./{rel_path}"},
        },
    }


def scan(spec_dir: Path) -> tuple[list[dict[str, Any]], list[str]]:
    manifests: list[dict[str, Any]] = []
    warnings: list[str] = []
    candidates = sorted(
        [*spec_dir.glob("*.yaml"), *spec_dir.glob("*.yml")]
    )
    if not candidates:
        warnings.append(f"no .yaml/.yml files found in {spec_dir}")
    for spec_path in candidates:
        try:
            spec = yaml.safe_load(spec_path.read_text(encoding="utf-8"))
        except yaml.YAMLError as exc:
            warnings.append(f"{spec_path.name}: YAML parse error: {exc}")
            continue
        if not isinstance(spec, dict) or "openapi" not in spec:
            warnings.append(f"{spec_path.name}: not an OpenAPI doc, skipped")
            continue
        manifests.append(manifest_for(spec_path, spec, spec_dir))
    return manifests, warnings


def main(argv: list[str] | None = None) -> int:
    parser = argparse.ArgumentParser(
        description="Emit Backstage-style API manifests for a spec folder."
    )
    parser.add_argument("spec_dir", type=Path, help="directory of OpenAPI YAML")
    parser.add_argument("--out", type=Path, default=None,
                        help="write manifests here (default: stdout)")
    args = parser.parse_args(argv)

    if not args.spec_dir.is_dir():
        print(f"error: not a directory: {args.spec_dir}", file=sys.stderr)
        return 2

    manifests, warnings = scan(args.spec_dir)
    for warning in warnings:
        print(f"warning: {warning}", file=sys.stderr)

    text = "---\n" + "---\n".join(
        yaml.safe_dump(m, sort_keys=False) for m in manifests
    )
    if args.out:
        args.out.write_text(text, encoding="utf-8")
        print(f"wrote {len(manifests)} manifest(s) to {args.out}")
    else:
        print(text, end="")
    return 0 if manifests else 1


if __name__ == "__main__":
    raise SystemExit(main())
\end{minted}

\begin{minted}{text}
---
apiVersion: backstage.io/v1alpha1
kind: API
metadata:
  name: northwind-robotics-fleet-api
  description: 'Manages the Northwind Robotics synthetic robot fleet: registration,
    health telemetry summaries, and decommissioning workflows.'
  annotations:
    northwind/spec-version: 2.4.0
spec:
  type: openapi
  lifecycle: published
  owner: team-fleet-platform
  definition:
    $text: ./specs/northwind_robots.yaml
---
apiVersion: backstage.io/v1alpha1
kind: API
metadata:
  name: northwind-robotics-telemetry-api
  description: Batch ingestion and querying of synthetic robot telemetry streams (motor
    current, battery voltage, wheel odometry).
  annotations:
    northwind/spec-version: 1.1.0
spec:
  type: openapi
  lifecycle: published
  owner: team-telemetry
  definition:
    $text: ./specs/northwind_telemetry.yaml
\end{minted}

\subsection{Listing 20.8 --- Smoke Tests for the Platform Tools}
\label{sec:ch20-listing-smoke}

A compact suite pins the cross-tool behaviors the lab relies on: docs
render classification labels, the funnel counts 60 developers and computes
the median TTFC of the fixture, malformed event lines are tolerated, the
tracker flags the shadow consumer and renders the T-0 sunset, and the
catalog scan emits one entry per spec.

\begin{minted}{python}
"""Smoke tests for the docs, funnel, deprecation, and catalog tools."""

from __future__ import annotations

from datetime import date
from pathlib import Path

import pytest
import yaml

import catalog_manifest
import deprecation_tracker
import docs_generator
import ttfc_funnel

HERE = Path(__file__).parent


def test_docs_generator_renders_endpoints_and_schemas() -> None:
    spec = yaml.safe_load(
        (HERE / "specs" / "northwind_robots.yaml").read_text(encoding="utf-8")
    )
    md = docs_generator.render_markdown(spec)
    assert "# Northwind Robotics Fleet API" in md
    assert "### `GET /robots`" in md
    assert "application/problem+json -> Problem" in md
    assert "[classification: internal]" in md  # data classification surfaced


def test_funnel_counts_and_ttfc() -> None:
    journeys, skipped = ttfc_funnel.load_journeys(
        HERE / "fixtures" / "onboarding_events.jsonl"
    )
    assert skipped == 0
    assert len(journeys) == 60
    assert len(ttfc_funnel.reached(journeys, "first_200")) == 24
    ttfc = ttfc_funnel.median_gap_minutes(journeys, "signup", "first_200")
    assert ttfc == pytest.approx(431.0)


def test_funnel_survives_malformed_lines(tmp_path: Path) -> None:
    bad = tmp_path / "events.jsonl"
    bad.write_text(
        '{"ts": "2026-03-02T09:00:00Z", "developer_id": "d1", '
        '"event": "signup"}\nnot json at all\n',
        encoding="utf-8",
    )
    journeys, skipped = ttfc_funnel.load_journeys(bad)
    assert skipped == 1 and len(journeys) == 1


def test_deprecation_tracker_blocked_and_schedule() -> None:
    registry = deprecation_tracker.load_registry(
        HERE / "fixtures" / "deprecation_registry.yaml"
    )
    consumers = deprecation_tracker.parse_consumers(registry)
    report = "\n".join(
        deprecation_tracker.readiness_report(
            registry, consumers, today=date(2026, 7, 15)
        )
    )
    assert "NOT READY TO RETIRE" in report
    assert "legacy-billing-hook" in report  # the shadow API shows up
    schedule = "\n".join(deprecation_tracker.schedule_report(registry))
    assert "T-0" in schedule and "brownout #2" in schedule


def test_catalog_manifest_scans_both_specs() -> None:
    manifests, warnings = catalog_manifest.scan(HERE / "specs")
    assert warnings == []
    names = {m["metadata"]["name"] for m in manifests}
    assert names == {
        "northwind-robotics-fleet-api",
        "northwind-robotics-telemetry-api",
    }
    assert all(m["kind"] == "API" for m in manifests)
    assert manifests[0]["spec"]["owner"].startswith("team-")
\end{minted}

Full suite:

\begin{minted}{text}
..............                                                           [100%]
14 passed in 0.59s
\end{minted}

%% ----------------------------------------------------------------------------
\section{Common Pitfalls \& Anti-Patterns}
\label{sec:ch20-pitfalls}

\subsection{Governance as Gatekeeping Theater}
\label{pit:theater}
A review board that meets fortnightly, demands slide decks, and rules on
naming by taste is not governance; it is theater with a queue. Symptoms:
review wait times measured in weeks, ``shadow reviews'' in hallway
conversations, teams shipping unreviewed APIs and begging forgiveness.
The fix is to move everything mechanical into lint (advisory first,
blocking later), cap review scope at judgment calls, and give the review
itself an SLA --- a gate slower than the bypass will be bypassed.

\subsection{Style Guide Without Linter Enforcement}
\label{pit:noenforce}
A wiki page of conventions decays the day it is published; enforcement
falls on senior reviewers, who become bottlenecks and then, inevitably,
inconsistent ones --- rule application becomes a function of which
reviewer you drew. If a rule cannot be expressed as executable data, the
honest options are to drop it or to reclassify it as advisory guidance.
Rules-as-data plus CI exit codes is the only enforcement mechanism that
scales past a handful of teams.

\subsection{The Portal Nobody Can Search}
\label{pit:portal}
A portal whose catalog is hand-curated decays into a museum: half the
entries are stale, search ranks marketing pages above reference docs, and
developers learn to ask in chat instead --- after which the portal team
declares ``developers don't use portals.'' Generate catalog entries from
the spec repo (Listing~20.7), render docs from the spec (Listing~20.4),
and instrument search queries with zero results; an unsearchable portal
is worse than none because it teaches developers to stop looking.

\subsection{SDKs Promised but Unmaintained}
\label{pit:sdk}
Announcing ``SDKs for seven languages'' at launch is a mortgage: every
spec change, generator bug, and language release comes due, in public,
under your brand. An unmaintained SDK with stale models teaches consumers
that your API cannot be trusted --- worse than no SDK, because the failure
arrives wearing your logo. Generate only what CI regenerates and
publishes automatically, and sunset SDKs with the same campaign
discipline as APIs.

\subsection{SLA Looser Than SLO, Reversed}
\label{pit:sla}
Promising a 99.95\% SLA while engineering to a 99.9\% SLO means the
contract runs out of budget before the team does: customers hit breach
before your alerting policy fires, and every incident becomes a legal
event. The invariant is SLA~$<$~SLO, with the gap sized as reaction time
in error-budget terms~\cite{beyer2016sre}. If sales demands SLA~$\geq$
SLO, the correct response is to raise the SLO \emph{and the engineering
investment}, not to sign the number.

\subsection{Sunsetting Without Telemetry}
\label{pit:notelemetry}
Announcing a kill date from a spreadsheet of ``known consumers'' retires
only the consumers you already knew about. The first brownout then
pages a team nobody mapped to an API key --- or worse, a production
workload whose owner left the company two reorganizations ago. Consumer
discovery \emph{must} start from gateway telemetry; the registry of
declared consumers is the escalation directory, not the census. No
census, no kill date.

\subsection{Docs Drifting From Spec}
\label{pit:docs-drift}
Hand-maintained reference docs diverge from the contract on the first
hotfix, and every divergence costs a consumer hours and costs you trust.
Docs-as-code (generate from OpenAPI, deploy from CI) makes drift a build
failure instead of a slow betrayal; the linter can even enforce that
every operation carries a \texttt{summary} so the generated docs are
never vacuous.

\begin{warningbox}
The compound anti-pattern is \textbf{governance without telemetry}:
style rules nobody measures, a portal nobody instruments, deprecations
announced to a census nobody verified. Each pitfall above is survivable
alone; together they produce the organization that discovers its real
API surface during an incident. If you fund only two things from this
chapter, fund the gateway consumer census and the TTFC funnel --- they
are the instruments every other control depends on.
\end{warningbox}

%% ----------------------------------------------------------------------------
\section{Hands-On Production Lab \& Real-World Case Study}
\label{sec:ch20-lab}

\subsection{Lab: Stand Up the Mini Platform}
Work from \texttt{code\textbackslash{}ch20\textbackslash{}} in PowerShell.
Expected duration: 60--90 minutes.

\begin{enumerate}
  \item \textbf{Install and baseline.} \texttt{pip install -r
        requirements.txt}, then \texttt{python -m pytest -q} --- all 14
        tests pass.
  \item \textbf{Lint the fleet.} Run \texttt{mini\_spectral.py} against
        both specs in \texttt{specs\textbackslash{}} and the
        non-compliant fixture; confirm exit codes 0, 0, 1 (use
        \texttt{\$LASTEXITCODE}). Identify which rule catches the most
        violations.
  \item \textbf{Author a rule.} Add a warn-severity rule
        \texttt{summary-required} to \texttt{ruleset.yaml} (hint: new rule
        kind not needed --- extend \texttt{check\_casing}-style iteration
        with a \texttt{require-operation-field} kind, or reuse
        \texttt{problem-json-errors} as a template). Verify the
        non-compliant fixture still exits 1 and the rule appears in
        output.
  \item \textbf{Generate docs.} Render the fleet spec with
        \texttt{docs\_generator.py --out fleet-api.md}; confirm the
        \texttt{x-classification} label appears in the Robot schema
        table.
  \item \textbf{Publish the catalog.} Run
        \texttt{catalog\_manifest.py .\textbackslash{}specs --out catalog-apis.yaml}
        and verify two entries with correct owners and lifecycle states.
  \item \textbf{Measure activation.} Run \texttt{ttfc\_funnel.py} on the
        fixture; write one sentence identifying the weakest funnel step
        and one DX investment that targets it.
  \item \textbf{Run the campaign.} Execute the deprecation tracker at
        \texttt{--today 2026-07-15} (expect exit 1, blocked) and again at
        \texttt{--today 2026-09-02} after editing all consumers to
        \texttt{migrated} (expect exit 0). Note how the communication
        schedule dates shift with the kill date.
\end{enumerate}

\subsection{Case Study: The Shadow API Nobody Owned}
\begin{examplebox}
\textbf{Northwind Robotics, incident INC-2291 (synthetic but typical).}
At 14:07 on a Tuesday, the fulfillment team deploys a routine change to
the \texttt{/shipments} service. Eight minutes later, checkout error
rates spike --- but checkout never calls \texttt{/shipments}. Forty
minutes of war-room archaeology later, gateway telemetry reveals the
coupling: an undocumented endpoint, \texttt{/shipments/internal-quote},
deployed eighteen months earlier by a contractor, serving 3{,}000 calls a
day from a billing integration whose owner field in the wiki pointed to a
team dissolved in the last reorganization. No catalog entry, no linter
history (the spec was never in the spec repo), no SLO, no on-call ---
and, fatally, no entry in the impact-assessment step of the fulfillment
team's change checklist, because as far as the platform was concerned the
endpoint did not exist.

The postmortem actions read like this chapter's table of contents:
(1)~diff observed gateway routes against the catalog weekly and quarantine
unregistered traffic behind an explicit registration step; (2)~require a
catalog entry --- and therefore an owner --- as an \emph{entry criterion
for} \textsf{published}; (3)~run the deprecation tracker's census step
against telemetry, not the wiki; (4)~treat ``unknown owner with nonzero
traffic'' as a severity-2 finding in its own right. The endpoint was
eventually owned, spec'd, linted, and published properly; the campaign to
retire it took four months and appears, disguised, as
\texttt{legacy-billing-hook} in Listing~20.6's registry.
\end{examplebox}

The lab's through-line is the loop of Figure~\ref{fig:ch20-loop}: lint and
catalog generation are the outward flow of governance; the funnel and the
campaign census are the telemetry flowing back. Running them side by side
makes the feedback tangible in under two hours.

%% ----------------------------------------------------------------------------
\section{Self-Assessment Exercises \& Deep-Dive Questions}
\label{sec:ch20-exercises}

\begin{enumerate}
  \item \textbf{State machine fluency.} For each event, name the
        transition and the evidence that must accompany it: (a)~a team
        opens a PR adding a new endpoint to a published spec;
        (b)~security flags an unauthenticated endpoint in a spec under
        review; (c)~a brownout reveals an unknown consumer three weeks
        before the kill date.
  \item \textbf{Rule authoring.} Write (in YAML, using the Listing~20.2
        rule kinds) a rule that requires every \texttt{POST} operation to
        document a \texttt{201} or \texttt{202} response. Which kind did
        you need, and what gap in the kind taxonomy did you find?
  \item \textbf{Severity economics.} Argue both sides: should
        \texttt{pagination-required-on-lists} be error or warn severity
        for a company at 40 APIs? At 400?
  \item \textbf{Funnel diagnosis.} The fixture's key$\rightarrow$first-call
        step converts at 73.3\% with a 194-minute median gap. Propose two
        distinct hypotheses and the instrument (not the fix) you would add
        to distinguish them.
  \item \textbf{SLA math.} An API measures 99.93\% monthly success. Its
        SLO is 99.9\%, its SLA 99.5\%. Compute both error budgets
        remaining. A change would spend another 0.05\%; which budget
        governs the decision and why?
  \item \textbf{Design the exception.} Draft the one-paragraph ADR for
        exempting a legacy partner endpoint from
        \texttt{problem-json-errors} for six months. What makes the
        exemption time-boxed and auditable rather than permanent?
  \item \textbf{Docs drift forensics.} Your generated docs and your
        hand-written quickstart disagree about a parameter name. Trace
        the failure: which artifact is wrong, which process allowed it,
        and what lint rule prevents recurrence?
  \item \textbf{Campaign design.} Using Listing~20.6's cadence, schedule
        a campaign for an API announced today with a 6-month notice
        period, one contracted enterprise consumer, and twelve internal
        consumers. Which cadence rows change for the enterprise consumer?
  \item \textbf{Shadow-API hunt.} Describe the exact telemetry query
        (fields, joins, thresholds) that surfaces an unregistered
        endpoint with real traffic, and the runbook for the first 24
        hours after it fires.
  \item \textbf{Platform ROI.} Your platform team of six costs \$1.4M/year.
        Construct the strongest honest one-slide ROI case using only
        adoption, TTFC, and incident-reduction evidence --- and identify
        the weakest assumption in your own slide.
\end{enumerate}

%% ----------------------------------------------------------------------------
\section{Interview Q\&A Vault}
\label{sec:ch20-interviews}

\begin{enumerate}[label=\textbf{Q\arabic*.}, leftmargin=1.7cm]
  \item \emph{[junior]} \textbf{Name the API lifecycle states and what
        fundamentally changes between \textsf{published} and
        \textsf{deprecated}.} Design, review, published (operate),
        deprecated, sunset, retired. At deprecation the contract freezes
        --- no new features, security fixes only --- the provider starts
        signaling (\texttt{Deprecation}/\texttt{Sunset} headers), and a
        dated migration campaign begins; the API still serves traffic
        under its SLA until the kill date.
  \item \emph{[junior]} \textbf{What belongs in an API style guide?}
        Only consumer-observable surfaces: naming, pagination, error
        shapes, authentication, versioning, idempotency, rate-limit
        signaling. Internal implementation choices stay with the team ---
        standardizing everything guarantees the guide is routed around.
  \item \emph{[junior]} \textbf{What is Spectral?} An open-source linter
        for OpenAPI: rulesets written as YAML data target parts of the
        spec (JSONPath-ish) with functions (casing, pattern, truthy,
        custom), run in CI so style-guide compliance is checked by
        machines, not reviewer memory.
  \item \emph{[junior]} \textbf{What is a developer portal?} The
        self-service front door: searchable API catalog, rendered docs,
        key issuance, sandbox access, status, and changelog. Its
        foundation is a machine-generated catalog of record; a
        hand-maintained portal decays into a museum.
  \item \emph{[junior]} \textbf{What is docs-as-code?} Documentation
        generated from the contract of record (OpenAPI), versioned and
        deployed by the same pipeline as the spec. Humans edit the spec;
        machines edit the docs --- so reference docs can never silently
        diverge from what the API does.
  \item \emph{[junior]} \textbf{What is TTFC and how do you reduce it?}
        Time-to-first-call: median elapsed time from signup to first
        successful (2xx) call. Reduce it by attacking the funnel step
        with the worst conversion: self-service instant keys
        (signup$\rightarrow$key), a 15-minute quickstart and sandbox
        (key$\rightarrow$call), and actionable error messages
        (call$\rightarrow$200). Instrument first; never optimize the
        funnel blind.
  \item \emph{[junior]} \textbf{SLA vs SLO vs SLI --- relationship?}
        SLI is the measurement, SLO the internal target on it, SLA the
        contractual promise with penalties. Invariant: SLA strictly
        looser than SLO, so the error-budget gap gives the team room to
        detect and fix before breaching the contract~\cite{beyer2016sre}.
  \item \emph{[junior]} \textbf{Deprecation vs sunset vs retirement?}
        Deprecation is the announcement (signal, freeze, campaign start);
        sunset is the enforcement window (brownouts, then permanent 410
        at the kill date); retirement is the aftermath (audit clean,
        catalog archived, runbooks removed).
  \item \emph{[mid]} \textbf{How do you enforce a style guide without
        killing velocity?} Convert rules to lint executed in CI
        (fast, impartial, self-service feedback in minutes), roll new
        rules out as warnings before errors, keep human review scoped to
        judgment with a two-day SLA, and invest the savings in the paved
        path so compliance is the easiest route --- enforcement should
        feel like autocomplete, not customs inspection.
  \item \emph{[mid]} \textbf{What makes a good custom lint rule?} A
        stable ID, an honest severity, a precise target, an actionable
        message with a JSON path to the offender, and a story --- the
        incident or inconsistency the rule exists to prevent. Rules
        without stories get litigated every review; rules with stories
        get quoted.
  \item \emph{[mid]} \textbf{When do generated SDKs hurt DX?} When the
        team cannot commit to owning them: generated code is unidiomatic
        in every language, lags language evolution, and a stale SDK
        under your brand is worse than none. Rule of thumb: generate when
        the spec is stable and CI regenerates and publishes
        automatically; otherwise ship an excellent OpenAPI doc, a
        quickstart, and a sandbox.
  \item \emph{[mid]} \textbf{Your funnel shows strong key issuance but
        weak first-call conversion. Diagnosis path?} The failure is
        between ``has credential'' and ``made a call'': quickstart
        quality, missing base-URL clarity, confusing scopes, or no
        sandbox. Instrument docs-viewed events and quickstart completion,
        time five developers through the getting-started flow, and check
        the \textsf{first call}$\rightarrow$\textsf{first 200} step too
        --- if 401s dominate first calls, the auth docs are the real
        culprit.
  \item \emph{[mid]} \textbf{What goes into a design-review packet?} The
        spec diff, the lint report, the resource model rationale, and an
        ADR per non-obvious decision (versioning strategy, pagination
        choice, error model deviations). The packet exists so the
        reviewer spends the two-day SLA window on judgment, not
        archaeology.
  \item \emph{[mid]} \textbf{Why ADRs for APIs specifically?} API
        decisions are unusually permanent --- consumers hard-code around
        them --- and unusually relitigated, because every new team
        re-asks ``why not GraphQL/why cursor pagination/why problem+json''.
        Recording rejected alternatives answers next year's mail without
        a meeting.
  \item \emph{[mid]} \textbf{How do you find out who consumes an API?}
        Gateway telemetry joined from API key/client ID to consumer
        identity --- ground truth, because packets don't lie. The catalog
        of record supplies owners and escalation contacts; the diff
        between observed and declared consumers is where shadow
        integrations live. Never run a census from a wiki.
  \item \emph{[mid]} \textbf{What is a paved path and how do you measure
        it?} The blessed, fully tooled route from idea to published API:
        repo template, pre-wired lint CI, generated docs, auto-catalog
        registration, sandbox, dashboards. Measure voluntary adoption
        (share of new APIs on the path) and time-to-publish versus the
        unpaved route; if teams need mandates, the path isn't paved well
        enough.
  \item \emph{[mid]} \textbf{Changelog discipline --- what does it look
        like operationally?} Every consumer-visible change, classified
        (added/changed/deprecated/removed/fixed/security), dated,
        shipped with the deploy, and machine-linked to deprecation
        headers. The changelog is a trust instrument: one lagging release
        teaches consumers to distrust every other signal you send.
  \item \emph{[mid]} \textbf{What makes a sandbox good?} Self-service
        instant keys, deterministic seeded synthetic data, no approvals,
        realistic failure modes (you can trigger a 429 and a 422 on
        purpose), and isolation from production. Its success metric is
        TTFC; a sandbox requiring a ticket is a waiting room with a URL.
  \item \emph{[senior]} \textbf{Design a sunset campaign for a v1 API
        with 200 consumers.} Census from gateway telemetry (expect
        unknowns); segment internal/external and by traffic; announce at
        T-180 via changelog, portal banner, email, and
        \texttt{Deprecation} headers; publish the kill date and a
        migration guide with a shim where feasible; reminders at T-90/60
        escalating to engineering leads; brownouts at T-30 (1 hour) and
        T-7 (4 hours) as escape hatches for non-readers; a documented,
        time-boxed extension process signed by the laggard's director;
        permanent 410 with a problem+json body at T-0; post-sunset
        traffic audit at T+7. Track readiness traffic-weighted and gate
        the retirement runbook on it.
  \item \emph{[senior]} \textbf{Why must the SLA be strictly looser than
        the SLO?} Because the gap between the internal budget
        ($1-\text{SLO}$) and the contractual budget ($1-\text{SLA}$) is
        reaction time measured in errors: alerting, rollback, and deploy
        freezes must consume SLO budget \emph{before} SLA budget is
        touched. SLA $\geq$ SLO means customers discover your incidents
        before your policy does, and every incident becomes a contract
        negotiation~\cite{beyer2016sre}.
  \item \emph{[senior]} \textbf{Brownouts: purpose, mechanics, risks?}
        Short, announced windows where the old version returns errors,
        flushing out consumers who ignore email before the permanent
        kill. Mechanics: gateway-level 410/404 with a problem body,
        per-consumer allowlist for carve-outs, business-hours timing so
        the right humans are awake. Risks: a brownout that pages the
        \emph{wrong} team (your census failed), and crying wolf ---
        brownouts that are always rolled back teach consumers to ignore
        them.
  \item \emph{[senior]} \textbf{Where do you draw the
        consistency-versus-autonomy line?} At the consumer-observable
        seam: everything a caller can see is standardized; everything
        behind the route table is the team's business. Revisit the line
        only with telemetry --- if an ungoverned internal choice keeps
        causing cross-team incidents, it earns promotion into the guide
        via ADR, not via decree.
  \item \emph{[senior]} \textbf{Kill date arrives; one consumer is
        unresponsive but still calling. What do you do?} Follow the
        pre-agreed policy, which existed precisely for this moment:
        confirm the escape-hatch criteria (is this a carve-out case or an
        abandonment case?), attempt final escalation through the census
        owner chain, brownout on schedule, and sunset --- with a
        documented reactivation path. Bending the kill date for silence
        teaches 199 other consumers that deadlines are negotiable.
  \item \emph{[senior]} \textbf{Design GDPR-style erasure propagation
        across an API ecosystem.} Erasure is an event, not a DELETE:
        a tombstone workflow fans out to the primary store, derived
        stores, caches, analytics, and downstream consumers you pushed
        data to, with per-hop acknowledgments; backups expire data by
        retention schedule rather than surgical edits. Schema-level
        classification labels tell you what to erase; the audit trail
        proves you did; quarterly rehearsals prove the workflow still
        works.
  \item \emph{[senior]} \textbf{How do you keep PII out of API logs at
        platform scale?} Don't rely on developer discipline: classify
        fields in the schema (\texttt{x-classification}), log only
        allow-listed metadata fields at the gateway and framework
        middleware level (never bodies by default), lint specs for
        classification coverage, and run periodic payload-sampling audits
        with automated detectors. The control lives in the pipeline, not
        in the style guide.
  \item \emph{[senior]} \textbf{Internal vs external deprecation policy
        --- what differs and why?} Notice period (sprint-plannable
        60--90 days vs contractual 6--12 months), enforcement lever (CI
        and deploy mandates vs persuasion, pricing, and brownouts within
        contract terms), and discovery (org chart available vs contracts
        and telemetry only). The substance is identical --- census,
        cadence, brownouts, kill date --- but external consumers hold
        contractual rights you must design around, not through.
  \item \emph{[senior]} \textbf{What telemetry should revise your style
        guide?} Lint violations by rule (a rule everyone suppresses is
        wrong or mistargeted), review findings that recur (recurrence
        means the guide is silent or unclear), support-ticket themes,
        funnel step conversions after docs changes, and error-mix shifts
        per consumer after releases. Governance without this feedback
        loop ossifies into theater.
  \item \emph{[staff]} \textbf{Build vs buy for portals and gateways ---
        your framework?} Buy the undifferentiated and mature (gateway,
        key management, hosted docs rendering); build what encodes your
        conventions (rulesets, paved-path templates, catalog enrichment
        glued to your org chart); default to composing: managed gateway,
        open-source portal core (Backstage-style), thin bespoke glue.
        The decisive question is not features but \emph{who operates it
        at 3 a.m.} and what it costs to leave.
  \item \emph{[staff]} \textbf{How should a platform team be funded, and
        why do cost-center models fail?} Fund as a product team with
        internal customers on a multi-year horizon, success measured by
        voluntary adoption. Cost-center funding turns the team into a
        ticket queue: roadmap by loudest stakeholder, no investment in
        self-service (which would ``reduce demand'' for the team), and
        eventual mandate-driven adoption that destroys the feedback loop
        that keeps the paved path paved~\cite{newman2021microservices}.
  \item \emph{[staff]} \textbf{How do you measure platform ROI?} Three
        columns: adoption (APIs on the paved path, traffic share behind
        the gateway), activation (TTFC median and funnel conversion,
        trended), and risk removed (lint-blocked non-compliant designs,
        incident rates of governed vs ungoverned APIs, sunsets completed
        without escalation). Report quarterly; translate to money only
        where the counterfactual is defensible.
  \item \emph{[staff]} \textbf{You inherit 500 APIs across 50 teams with
        zero governance. First 90 days?} Days 1--30: instrument ---
        gateway census, catalog skeleton generated from whatever specs
        exist, TTFC baseline. Days 31--60: publish the v1 style guide
        (consumer-observable surfaces only), ship the linter as
        \emph{advisory}, and stand up one paved path for \emph{new} APIs.
        Days 61--90: flip two high-value rules to blocking, run the first
        review-SLA'd design reviews, and publish the first transparency
        report. Never begin with a mandate to retrofit all 500.
  \item \emph{[staff]} \textbf{A strong team demands an exception to the
        error-shape rule ``just this once.'' Your process?} Exceptions
        are ADRs, not conversations: written context, explicit
        time-boxing, a named owner, a suppression comment in the spec
        referencing the ADR, and a review date. Permanent undocumented
        exceptions compound; time-boxed documented ones are governance
        working as intended --- the rule learns where its edges are.
  \item \emph{[staff]} \textbf{Your largest revenue account refuses to
        migrate off v1. Sunset anyway?} This is a business decision the
        campaign should have surfaced at T-90, not T-7. Options, in
        order: dedicated migration engineering (often cheaper than the
        standoff), a paid extended-support tier with a contractual final
        date, or accepting v1 as a frozen, ring-fenced deployment with
        its own cost model charged back. What you don't do is quietly
        extend forever --- an undead v1's maintenance cost is invisible
        precisely because nobody is measuring it.
  \item \emph{[staff]} \textbf{What organizational signals tell you
        governance has become theater?} Review queues measured in weeks
        while ship rates stay flat (teams are bypassing); lint suppressions
        clustering on the same rules; a catalog whose entries disagree
        with gateway telemetry; style-guide pages last edited by someone
        who left; and --- the tell --- governance metrics that report
        \emph{activity} (reviews held) instead of \emph{outcomes}
        (consistency, adoption, TTFC, sunset success rate).
\end{enumerate}

%% ----------------------------------------------------------------------------
\section{External Bibliography \& Further Reading}
\label{sec:ch20-bibliography}

\begin{itemize}
  \item \textbf{OpenAPI Specification v3.1.0}~\cite{openapi31} --- the
        contract of record this chapter's linter, docs generator, and
        catalog all consume.
  \item \textbf{Gehl \& Higginbotham, \emph{API Design Patterns}
        (Manning, 2021)}~\cite{gehl2021apidesign} --- pattern-level
        grounding for the style-guide surfaces: naming, pagination,
        versioning, and long-running operations.
  \item \textbf{Beyer et al., \emph{Site Reliability Engineering}
        (O'Reilly, 2016)}~\cite{beyer2016sre} --- SLOs, error budgets,
        and the SLA~$<$~SLO doctrine.
  \item \textbf{Newman, \emph{Building Microservices}, 2nd ed.
        (O'Reilly, 2021)}~\cite{newman2021microservices} --- platform and
        team-boundary thinking behind paved paths.
  \item \textbf{The Twelve-Factor App}~\cite{twelvefactor} ---
        config-in-environment discipline that sandbox and portal key
        issuance builds on.
  \item \textbf{Spectral documentation},
        \texttt{https://docs.stoplight.io/docs/spectral} --- ruleset
        anatomy, built-in functions, custom rules; the production
        counterpart of Listing~20.2.
  \item \textbf{Backstage documentation},
        \texttt{https://backstage.io/docs} --- software catalog and
        \texttt{kind: API} entity model behind Listing~20.7.
  \item \textbf{openapi-generator},
        \texttt{https://openapi-generator.tech} --- the reference SDK
        generation pipeline, including its templating and CI regeneration
        story.
  \item \textbf{Redocly / Redoc} (\texttt{https://redocly.com/docs}) and
        \textbf{Stoplight Elements} --- production docs-as-code renderers
        whose role Listing~20.4 miniaturizes.
  \item \textbf{Skelton \& Pais, \emph{Team Topologies} (IT Revolution,
        2019)} --- the platform-team-as-product funding and interaction
        models of Section~\ref{sec:ch20-economics}.
  \item \textbf{Google AIP program}, \texttt{https://google.aip.dev} ---
        a mature, public, lint-enforced API design standard; study its
        AIP-122 (resource names) and AIP-158 (pagination) alongside
        Table~\ref{tab:ch20-styleguide}.
  \item \textbf{Stripe, Slack, and Twilio engineering blogs}
        (\texttt{https://stripe.com/blog}, \texttt{https://slack.engineering},
        \texttt{https://www.twilio.com/blog}) --- primary sources on DX
        engineering at scale: Stripe's API versioning and changelog
        discipline, Slack's developer-platform iterations, Twilio's
        helper-library (SDK) maintenance economics.
  \item \textbf{RFC~8594 (the \texttt{Sunset} header)} and the
        deprecation-header draft --- the wire-level machinery behind
        Section~\ref{sec:ch20-deprecation}, introduced in
        Chapter~\ref{ch:versioning}.
\end{itemize}

%% ----------------------------------------------------------------------------
\section{Capstone Projects}
\label{sec:ch20-capstones}

\subsection*{Capstone 20.1 [junior] --- Three New Rules for the Fleet}
\textbf{Objective:} Extend the Northwind ruleset with three custom rules
and prove each one bites.

\textbf{Inputs/Datasets:} \texttt{code\textbackslash{}ch20\textbackslash{}}
(\texttt{mini\_spectral.py}, \texttt{ruleset.yaml}, both spec fixtures).

\textbf{Steps:}
\begin{enumerate}
  \item Add rules: \texttt{summary-required} (every operation has a
        non-empty \texttt{summary}), \texttt{no-verbs-in-paths}
        (path segments match a noun-only denylist: \texttt{get},
        \texttt{create}, \texttt{delete}, \texttt{fetch}),
        \texttt{cursor-type-string} (any \texttt{cursor} parameter is
        \texttt{type: string}).
  \item Extend the linter with any new rule kinds needed, keeping rules
        as data.
  \item Add one violating case per rule to a copy of the non-compliant
        fixture; assert exit codes and rule IDs in new tests.
\end{enumerate}

\textbf{Acceptance Criteria:}
\begin{itemize}
  \item[$\square$] Compliant specs still exit 0; the extended fixture
        exits 1 with all three new rule IDs.
  \item[$\square$] Each violation carries a JSON path pinpointing the
        offending node.
  \item[$\square$] New rule kinds are unit-tested, including a
        bad-ruleset (unknown kind) test exiting 2.
\end{itemize}

\textbf{Stretch Goals:} Add per-rule suppression comments
(\texttt{x-lint-ignore: [rule-id]}) honored by the engine and logged as
warnings.

\subsection*{Capstone 20.2 [mid] --- Docs-as-Code Pipeline}
\textbf{Objective:} Assemble lint $\rightarrow$ docs $\rightarrow$
catalog into one offline pipeline script with a CI gate.

\textbf{Inputs/Datasets:} \texttt{code\textbackslash{}ch20\textbackslash{}};
the two specs under \texttt{specs\textbackslash{}}.

\textbf{Steps:}
\begin{enumerate}
  \item Write \texttt{pipeline.py}: lint every spec (block on errors),
        render Markdown docs to \texttt{docs\textbackslash{}}, and emit
        \texttt{catalog-apis.yaml}.
  \item Add a staleness check: if a spec changed but its rendered docs
        did not regenerate (compare content hashes stored in a manifest),
        exit 1 with a clear message.
  \item Record a one-page ADR describing the pipeline's blocking policy.
\end{enumerate}

\textbf{Acceptance Criteria:}
\begin{itemize}
  \item[$\square$] One command produces docs and catalog for all specs;
        a non-compliant spec blocks the run.
  \item[$\square$] Editing a spec without rerunning the generator fails
        the staleness gate.
  \item[$\square$] Everything runs offline; output is byte-deterministic
        across runs.
\end{itemize}

\textbf{Stretch Goals:} Add a Markdown link checker (stdlib) over the
generated docs.

\subsection*{Capstone 20.3 [mid] --- TTFC Instrumentation Pack}
\textbf{Objective:} Extend the funnel analyzer from report to decision
tool.

\textbf{Inputs/Datasets:} \texttt{ttfc\_funnel.py},
\texttt{fixtures\textbackslash{}onboarding\_events.jsonl} (222 synthetic
events, 60 developers).

\textbf{Steps:}
\begin{enumerate}
  \item Add a \texttt{--by channel} breakdown (the fixture's
        \texttt{channel} field) computing per-channel funnels.
  \item Add a weekly cohort table (signups grouped by ISO week) with
        per-cohort TTFC medians.
  \item Write a ``DX experiment'' report: propose one intervention,
        state which funnel step it targets, and the metric movement that
        would confirm it.
\end{enumerate}

\textbf{Acceptance Criteria:}
\begin{itemize}
  \item[$\square$] Per-channel funnels show conversion deltas with
        counts (no percentage without an $n$).
  \item[$\square$] Cohort table is deterministic and sorted; malformed
        events are still counted and skipped, never fatal.
  \item[$\square$] The experiment report names one metric, one step, and
        one falsifiable threshold.
\end{itemize}

\textbf{Stretch Goals:} Flag small-$n$ cohorts ($n < 10$) in the output
to discourage over-reading noise.

\subsection*{Capstone 20.4 [senior] --- Sunset Campaign for Fleet API v1}
\textbf{Objective:} Design and simulate the full retirement of a legacy
version with heterogeneous consumers.

\textbf{Inputs/Datasets:} \texttt{deprecation\_tracker.py},
\texttt{fixtures\textbackslash{}deprecation\_registry.yaml}; a 12-consumer
registry you extend it into (add tiers, contract dates, one shadow
consumer).

\textbf{Steps:}
\begin{enumerate}
  \item Extend the registry schema with \texttt{notice\_days} per tier
        and \texttt{extension\_until} per consumer; validate
        defensively.
  \item Add a simulation mode: given a migration-velocity assumption per
        tier, project residual traffic weekly until the kill date.
  \item Produce the campaign plan: cadence rows adjusted per tier,
        brownout scope, carve-out list, and the go/no-go checklist for
        T-0.
\end{enumerate}

\textbf{Acceptance Criteria:}
\begin{itemize}
  \item[$\square$] The projection flags, in advance, the consumers who
        will miss the date under stated velocity.
  \item[$\square$] The plan differentiates internal vs external policy
        per Table~\ref{tab:ch20-internal-external}.
  \item[$\square$] Exit codes remain CI-meaningful; all dates are
        explicit CLI inputs (deterministic).
\end{itemize}

\textbf{Stretch Goals:} Model one extension request and show its effect
on the brownout schedule for non-extended consumers.

\subsection*{Capstone 20.5 [staff] --- Platform ROI Scorecard \& Governance Operating Model}
\textbf{Objective:} Design the quarterly platform review for a fictional
50-team Northwind Robotics.

\textbf{Inputs/Datasets:} This chapter's tools plus invented-but-plausible
baselines (document all assumptions); no real company data.

\textbf{Steps:}
\begin{enumerate}
  \item Define the scorecard: adoption, activation (TTFC funnel), risk
        removed; a data source and an owner per metric.
  \item Build \texttt{scorecard.py}: reads the funnel fixture, a lint
        history CSV you synthesize, and an incident log you synthesize;
        renders the one-page scorecard.
  \item Write the operating model: rule lifecycle (propose $\rightarrow$
        advisory $\rightarrow$ blocking $\rightarrow$ retire), review
        SLA, exception process, and the platform team's funding ask with
        the ROI narrative.
\end{enumerate}

\textbf{Acceptance Criteria:}
\begin{itemize}
  \item[$\square$] Every scorecard number traces to a reproducible
        computation from synthetic inputs.
  \item[$\square$] The ROI narrative states its weakest assumption
        explicitly.
  \item[$\square$] The operating model fits on two pages and names
        decision rights for every governance event.
\end{itemize}

\textbf{Stretch Goals:} Add a counterfactual section: estimated incident
cost avoided using governed-vs-ungoverned incident rates, with the
selection-bias caveat addressed.

%% ----------------------------------------------------------------------------
\section{Python Dependencies Appendix}
\label{sec:ch20-dependencies}

The chapter's code targets Python~3.11+ and needs four third-party
packages (only PyYAML and pytest are exercised at runtime; the other two
support the program-level patterns referenced below):

\begin{minted}{text}
pyyaml>=6.0
jsonschema>=4.21
pydantic>=2.7
pytest>=8
\end{minted}

Install on Windows/PowerShell (virtual environment recommended):

\begin{minted}{powershell}
cd code\ch20
python -m venv .venv
.\.venv\Scripts\Activate.ps1
pip install -r requirements.txt
\end{minted}

\subsection{PyYAML (\texttt{pyyaml>=6.0})}
\textbf{Description:} The de-facto YAML parser for Python.
\textbf{Why included:} OpenAPI documents and rulesets in this chapter are
YAML; \texttt{yaml.safe\_load} parses specs, rulesets, registries, and
catalog emission uses \texttt{yaml.safe\_dump}. Safe loaders only ---
\texttt{yaml.load} without a safe loader is a code-execution hazard on
untrusted input.
\textbf{Alternatives:} \texttt{ruamel.yaml} (round-trip editing with
comments preserved, useful if you want a linter that \emph{fixes} specs);
the real Spectral CLI (Node.js) for production-grade linting; Redocly CLI
for lint-plus-bundle.
\textbf{Installation:} \texttt{pip install pyyaml>=6.0} (PowerShell;
quote as \texttt{"pyyaml>=6.0"} if your shell interprets \texttt{>}).

\subsection{jsonschema (\texttt{jsonschema>=4.21})}
\textbf{Description:} JSON Schema validation for Python; OpenAPI~3.1
schemas are JSON Schema dialects~\cite{openapi31}.
\textbf{Why included:} The natural next layer after structural linting is
\emph{schema-aware} rules --- e.g., validating that a spec's example
payloads actually satisfy their declared schemas, or that a registry YAML
conforms to a published registry schema (the defensive parsing in
Listings~20.2 and~20.6 is the hand-rolled floor; \texttt{jsonschema} is
the ceiling).
\textbf{Alternatives:} \texttt{openapi-spec-validator} and
\texttt{openapi-schema-validator} (OpenAPI-aware validation);
Spectral/Redocly for schema checks embedded in a lint pipeline.
\textbf{Installation:} \texttt{pip install "jsonschema>=4.21"}.

\subsection{Pydantic (\texttt{pydantic>=2.7})}
\textbf{Description:} Type-annotated data modeling and validation.
\textbf{Why included:} The chapter's tools deliberately use stdlib
dataclasses for transparency, but a production platform models its
catalog entries, registry records, and ruleset schema as pydantic models
--- gaining validation errors with field paths, JSON Schema export for
free, and parity with the FastAPI services of
Chapter~\ref{ch:fastapi}~\cite{openapi31}.
\textbf{Alternatives:} stdlib \texttt{dataclasses} (as used here) plus
\texttt{jsonschema} for boundary validation; \texttt{msgspec} for
speed-critical parsing; \texttt{attrs} for lightweight class ergonomics.
\textbf{Installation:} \texttt{pip install "pydantic>=2.7"}.

\subsection{pytest (\texttt{pytest>=8})}
\textbf{Description:} The standard Python test runner.
\textbf{Why included:} Both test suites (Listings~20.3 and~20.8) use
pytest's \texttt{tmp\_path} and \texttt{capsys} fixtures to exercise CLI
exit codes --- the CI contract of the linter and tracker --- without
shelling out, keeping tests offline and deterministic per
Chapter~\ref{ch:testing}.
\textbf{Alternatives:} \texttt{unittest} (stdlib, more ceremony);
\texttt{nox}/\texttt{tox} as session runners on top of pytest; MkDocs is
\emph{not} a test alternative but is the natural companion for publishing
this chapter's generated Markdown as a docs site.
\textbf{Installation:} \texttt{pip install "pytest>=8"}.

%% ----------------------------------------------------------------------------
\section{Chapter Summary \& Key Takeaways}
\label{sec:ch20-summary}

\begin{itemize}
  \item \textbf{The lifecycle is a state machine, not a poster.} Seven
        states with evidence-based entry/exit criteria make ownership,
        tooling, and audits mechanical. The \textsf{deprecated} state is
        first-class: frozen contract, \texttt{Sunset} headers, live
        campaign.
  \item \textbf{Governance compiles or it doesn't exist.} Style guides
        standardize only the consumer-observable seam; Spectral-style
        rulesets enforce them as data in CI with exit codes that
        distinguish ``non-compliant'' from ``gate broken''; review keeps
        a two-day SLA and advisory-by-default verdicts; ADRs make
        exceptions and rejected alternatives permanent public record.
  \item \textbf{DX is a measured funnel.} Portals and catalogs are
        generated, not curated; docs are build artifacts of the spec;
        SDKs are generated only when CI regenerates them forever; TTFC
        (signup $\rightarrow$ key $\rightarrow$ call $\rightarrow$ 200)
        is the activation metric, and every DX investment names the
        funnel step it moves.
  \item \textbf{APIs are products with contract discipline.} Tier on a
        value metric enforceable at the gateway; keep SLA strictly
        looser than SLO so error budget is reaction
        time~\cite{beyer2016sre}; changelogs are trust instruments.
  \item \textbf{Sunsetting is a campaign run on telemetry.} The census
        comes from the gateway, not the wiki; cadence runs T-180 to T+7
        with brownouts as escape hatches; internal and external consumers
        get different levers but the same dates.
  \item \textbf{Data governance rides the contract.} Classification
        labels in schemas, payload minimization, redaction by default,
        retention clocks, and rehearsed deletion propagation across
        stores, caches, and downstream consumers.
  \item \textbf{Platform economics justify the platform.} Buy the
        undifferentiated, build the conventional; fund the platform team
        as a product team~\cite{newman2021microservices}; report ROI in
        adoption, TTFC, and risk removed --- or watch the budget leave.
\end{itemize}

%% ----------------------------------------------------------------------------
\section{Quick-Reference Card}
\label{sec:ch20-quickref}

\begin{table}[htbp]
\centering\small
\begin{tabularx}{\textwidth}{l X X}
\toprule
\textbf{State} & \textbf{Consumer may assume} & \textbf{Provider must do} \\
\midrule
design & nothing & spec in repo, ADR for hard choices \\
review & nothing & lint green; review within SLA \\
published & stable contract, SLO, docs & operate, changelog, compatible evolution only \\
deprecated & works until kill date; no new features & headers on, campaign running, security fixes only \\
sunset & brownouts then 410 & final warnings, 410 + problem body with migration link \\
retired & nothing (404/410 forever) & audit clean, catalog archived \\
\bottomrule
\end{tabularx}
\caption{Lifecycle states: the contract each side is under.}
\label{tab:ch20-states-quick}
\end{table}

\begin{table}[htbp]
\centering\small
\begin{tabularx}{\textwidth}{l X}
\toprule
\textbf{Surface} & \textbf{Core rule (lint-enforced)} \\
\midrule
Naming & kebab-case path segments; camelCase \texttt{operationId} \\
Versioning & \texttt{info.version} strict semver; URI major version \\
Errors & every operation declares \texttt{application/problem+json} on 4xx/5xx/default \\
Pagination & collection GETs accept \texttt{limit} + \texttt{cursor} \\
Ownership & \texttt{info.contact} + \texttt{x-lifecycle.owner} required \\
Governance hygiene & rules as data; warn $\rightarrow$ error rollout; suppressions are ADR-linked \\
\bottomrule
\end{tabularx}
\caption{Style-guide core rules and their enforcement.}
\label{tab:ch20-rules-quick}
\end{table}

\begin{table}[htbp]
\centering\small
\begin{tabularx}{\textwidth}{l X}
\toprule
\textbf{Metric} & \textbf{Definition / target practice} \\
\midrule
TTFC & median $t(\text{first 200}) - t(\text{signup})$; the activation metric \\
Funnel conversion & $C_i = N_i/N_{i-1}$ per step; fix the weakest step first \\
Census coverage & \% of gateway traffic mapped to cataloged consumers; target 100\%, diff weekly \\
Paved-path adoption & \% of new APIs via the paved path; voluntary, trended \\
SLA/SLO gap & SLA $<$ SLO; gap = reaction budget in errors~\cite{beyer2016sre} \\
Campaign readiness & residual traffic on non-migrated consumers; T-0 gate \\
Portal health & zero-result searches, docs-to-spec staleness hash \\
\bottomrule
\end{tabularx}
\caption{DX and governance metrics cheat sheet.}
\label{tab:ch20-metrics-quick}
\end{table}

\section{Standardized Evidence Addendum}
\subsection{Benchmark Snapshot Template}
\begin{table}[h]
  \centering
  \small
  \begin{tabularx}{\textwidth}{l c c c c}
    \toprule
    \textbf{Config} & \textbf{p50 latency} & \textbf{p99 latency} & \textbf{Error rate} & \textbf{Throughput} \\
    \midrule
    Baseline & -- & -- & -- & 1.00x \\
    Candidate A & -- & -- & -- & -- \\
    Candidate B & -- & -- & -- & -- \\
    \bottomrule
  \end{tabularx}
  \caption{Minimum benchmark table to keep per chapter.}
\end{table}

\subsection{Request Trace Template}
\begin{itemize}
  \item \textbf{Request:} one representative API call for this chapter's topic.
  \item \textbf{Before:} behavior (headers, payload, latency, failure mode) with the naive approach.
  \item \textbf{After:} behavior with the technique in this chapter.
  \item \textbf{Result:} what changed in correctness, resilience, latency, and operability.
\end{itemize}

\subsection{Cost and Latency Budget Snapshot}
\begin{table}[h]
  \centering
  \small
  \begin{tabularx}{\textwidth}{l c c}
    \toprule
    \textbf{Stage} & \textbf{Latency budget} & \textbf{Cost driver} \\
    \midrule
    TLS / connection setup & -- & handshake round trips, keep-alive hit rate \\
    Request validation & -- & schema complexity, CPU per request \\
    Business logic and I/O & -- & downstream fan-out, query cost \\
    Serialization and egress & -- & payload size, compression ratio \\
    \bottomrule
  \end{tabularx}
  \caption{Operational budget template for this chapter's technique.}
\end{table}

\section{Configuration Preset Template}
\begin{minted}{text}
# api_policy.yaml
policy_version: "v1.0.0"
component: "<chapter_topic>"
timeout_ms: 0
retry_budget: 0
fallback_strategy: "fail_fast"
rollback_trigger: "error_budget_burn_gt_threshold"
\end{minted}

\section{CI Quality Gate Specification}
\begin{itemize}
  \item contract tests green against the pinned OpenAPI/AsyncAPI/Protobuf spec,
  \item no breaking schema changes without a version bump,
  \item error responses conform to RFC 9457 problem-details shape,
  \item benchmark deltas within approved regression bounds (latency and throughput),
  \item policy version and rollback plan present in release metadata.
\end{itemize}

\section{Incident Runbook Template}
\begin{enumerate}
  \item Detect regression from SLO burn-rate alerts or consumer reports.
  \item Scope impacted endpoints, versions, and consumer segments.
  \item Contain impact (rollback, feature flag, rate-limit tightening, or traffic shift).
  \item Diagnose with traces, structured logs, and contract diffs.
  \item Recover with validated fix and verified rollout procedure.
  \item Prevent recurrence via new regression tests and CI gates.
\end{enumerate}

\section{Cross-Chapter Decision Card}
\begin{table}[h]
  \centering
  \small
  \begin{tabularx}{\textwidth}{l X X}
    \toprule
    \textbf{Dimension} & \textbf{Choose this approach when} & \textbf{Prefer an alternative when} \\
    \midrule
    Use case & -- & -- \\
    Latency budget & -- & -- \\
    Operational cost & -- & -- \\
    Team maturity & -- & -- \\
    \bottomrule
  \end{tabularx}
  \caption{Decision card to complete per chapter.}
\end{table}

\section{ROI Model and Build-vs-Buy}
\begin{itemize}
  \item \textbf{Build when:} the capability is differentiating, compliance forbids vendors, or scale makes vendor pricing irrational.
  \item \textbf{Buy when:} the capability is commodity (auth, gateways, observability), time-to-market dominates, or staffing is thin.
  \item \textbf{Hidden costs to model:} on-call load, upgrade cadence, expertise ramp, data egress fees.
\end{itemize}

\section{Disaster Recovery Notes (RPO/RTO)}
\begin{itemize}
  \item Define recovery point objective (RPO): maximum acceptable data loss window.
  \item Define recovery time objective (RTO): maximum acceptable restoration time.
  \item Test restores regularly; an untested backup is not a backup.
\end{itemize}

